\documentclass[11pt]{article}
\usepackage{amsmath,amssymb,amsthm}
\usepackage{array,longtable}
\usepackage[margin=1in]{geometry}
\usepackage{xurl}
\usepackage[hidelinks]{hyperref}
\setcounter{tocdepth}{1}
\newtheorem{lemma}{Lemma}
\title{Notes on two prime power partitions with the same bases}
\author{}
\date{}
\begin{document}
\maketitle
% Source map: the abstract and opening setup retain diophantine_analysis.tex.
% The quadratic-form proofs and elementary reductions consolidate that source.
% Valuations, size bounds, orders, and the common-gap construction consolidate
% even_odd_case.tex. Later additions are marked at their source locations below.
\section*{Abstract}

Let $p > 5$ be a prime. We ask how many solutions exist where $p = 2^m + q^n$ for $m, n \geq 1$ and $q$ a fixed prime.
We know of a single example where $p$ has two solutions, $p =  11$. Do others exist? First we proceed by
considering basic arithmetic obstructions that prevent certain solutions. This breaks down to two main cases,
based on whether a common factor of putative solutions, $A$, is either greater than or equal to 1.

\section*{Earlier proof route}
\addcontentsline{toc}{section}{Earlier proof route}

The classification is proved in Section~\ref{sec:quadratic-norm}.
Section~\ref{sec:proof-reading} lists its prerequisites. The argument
below records the route through orders and reciprocity.

Write $M=m_1-m_2$ and $N=n_2-n_1$. Subtracting the partitions
gives $Aq^{n_1}=2^M-1$ and $A2^{m_2}=q^N-1$, with $\gcd(A,2q)=1$.
The quadratic form gives $M$ even and $N$ odd. Factoring $2^M-1$
then gives $p=11$ when $A=1$. For $q=3$, any $n_1>1$ forces $7\mid A$,
but $\operatorname{ord}_7(3)=6$. The case $n_1=1$ again gives $p=11$
(Lemmas~\ref{lem:Aone} and~\ref{lem:qthree}).

For $A>1$ and $q>3$, the parity and congruence arguments give
$q\equiv1\pmod3$, $m_2=v_2(q-1)$ even, $n_1$ even, and
$N\ge5$ with $\gcd(N,6)=1$. Put $r=\operatorname{ord}_q(2)$ and
$\rho=v_q(2^r-1)$. Lifting gives $n_1=\rho+v_q(M)$.

The two large powers yield
$0<n_2\log q-m_1\log 2<(27/26)q^{-N}$.
The fixed-$q$ estimate bounds $M$. Keeping the two logarithmic heights
separate also gives the uniform bound $N\le2519$ in the remaining
case (Lemma~\ref{lem:uniform-N}). This leaves finitely many
dyadic approximation problems; Roth's theorem gives finiteness
over all $q$. For bounded $m_2$, Schinzel--Tijdeman gives an
effective bound on $q$.
Lemma~\ref{lem:odd-gap-part} excludes $M=2^a$ and $M=3\cdot2^a$
for every $a$ by computing the odd part of $q-1$ modulo $1024$.
If $q\mid M$, then $qr\mid M$, giving absolute bounds $q<10^6$
and $M<4\cdot10^6$. The finite check in
Section~\ref{sec:q-divides-exclusion} excludes this branch.

It remains to consider $q\nmid M$ and $n_1=\rho\ge2$ even.
The modulo-$5$ argument gives $q\equiv1\pmod{48}$
(Lemma~\ref{lem:q-one-mod-eight}). The size and residue arguments
in Section~\ref{sec:additional-returns} give $7\nmid\gcd(M,N)$ and $\gcd(M,N)=1$ for $N<65$.
With $E=\operatorname{lcm}(2,r)$, we have
\[
E\mid M,\qquad \gcd(E,N)=1,\qquad
a=\frac{2^E-1}{q^\rho}\mid A,\qquad q^N\equiv1\pmod a.
\]
Thus one prime-power factor of $a$ at which $q$ has order sharing
a prime divisor with $E$ would exclude every $N$ for that $q$.
Section~\ref{sec:reciprocity} computes products of power-residue
symbols that can certify such a factor without factoring $a$.
The open step on this route is to force an order obstruction for every
remaining $q$.

\clearpage
{\small\tableofcontents}
\clearpage
\section{Two partitions and the common factor}
\label{sec:setup}

Let $p$ be prime with two prime power solutions sharing the same base $q$. So, we have
\[
p = 2^{m_1} + q^{n_1} = 2^{m_2} + q^{n_2}, \qquad m_1 > m_2, \quad n_2 > n_1.
\]
So then writing $\Delta m = m_1 - m_2, \Delta n = n_2 - n_1$, then just following our noses we
wind up with a coprime quantity to our bases, which is our central obstacle. That is, with

\[
2^{m_2}\left(2^{\Delta m} - 1\right) = q^{n_1}\left(q^{\Delta n} - 1\right).
\]
and  $\gcd(2,q) = 1$, we have $2^{m_2} \mid q^{\Delta n} - 1$ and $q^{n_1} \mid 2^{\Delta m} - 1$.

And so we define
\[
A = \frac{2^{\Delta m} - 1}{q^{n_1}} = \frac{q^{\Delta n} - 1}{2^{m_2}},
\]
Now we consider
\[
A\,q^{n_1} = 2^{\Delta m} - 1, \qquad A\,2^{m_2} = q^{\Delta n} - 1, \qquad A < 2^{\Delta m}, \quad A < q^{\Delta n}.
\]

The coprimality of $A$ with the bases follows directly:
$Aq^{n_1}=2^{\Delta m}-1$ makes $A$ odd, and
$A2^{m_2}=q^{\Delta n}-1$ gives $q\nmid A$. Thus
\[
\gcd(A,2q)=1.
\]
In particular $A$ is a positive integer, so the two cases are $A=1$ and
$A>1$.

\subsection{The quadratic form and gap parity}

Using $2^{\Delta m}=Aq^{n_1}+1$, we can write
\[
p=A2^{m_2}q^{n_1}+2^{m_2}+q^{n_1}.
\]
Consequently
\begin{align*}
Ap+1
 &= (A2^{m_2}+1)(Aq^{n_1}+1)\\
 &= q^{\Delta n}2^{\Delta m},
\end{align*}
so
\[
\boxed{Ap=2^{\Delta m}q^{\Delta n}-1.}
\]
Write
\[
\Delta m=2k+e_2,\qquad \Delta n=2j+e_q,
\qquad e_2,e_q\in\{0,1\},\quad k,j\geq0,
\]
and set
\[
D=2^{e_2}q^{e_q},\qquad X=2^kq^j.
\]
Then $D$ is squarefree and
\[
\boxed{Ap=DX^2-1.}
\]
Since $p>2,q$, reduction modulo $p$ gives
\[
DX^2\equiv1\pmod p,\qquad
\left(\frac{D}{p}\right)=1.
\]
The two defining equations for $A$ also give
\begin{equation}\label{eq:A-size}
A^2
 =\frac{(2^{\Delta m}-1)(q^{\Delta n}-1)}{2^{m_2}q^{n_1}}
 <\frac{2^{\Delta m}q^{\Delta n}}{2^{m_2}q^{n_1}}
 \leq\frac{DX^2}{6}.
\end{equation}
Here $2^{m_2}q^{n_1}\geq6$.

\begin{lemma}\label{lem:squares}
The gaps $\Delta m$ and $\Delta n$ cannot both be even.
\end{lemma}
\begin{proof}
If both are even, $D=1$ and $X=2^kq^j\geq6$. Thus
\[
Ap=(X-1)(X+1),\qquad A<X/\sqrt6<X-1.
\]
Since $p$ is prime, it divides at least one of $X-1,X+1$.
The other factor then divides $A$, which is impossible because both
factors exceed $A$.
\end{proof}

\begin{lemma}\label{lem:parity}
The gap $\Delta m$ is even. Consequently $\Delta n$ is odd.
\end{lemma}
\begin{proof}
Since $p>5$, it is nonzero modulo $3$. Exactly one parity of $m$
satisfies $2^m\equiv p\pmod3$.
If $q=3$, both $m_1,m_2$ have that parity.
If $q\neq3$, neither has that parity, since $p-2^{m_i}$ is a power
of $q$ and is therefore nonzero modulo $3$.
In either case $m_1,m_2$ have the same parity, so $\Delta m$ is even.
Lemma~\ref{lem:squares} now makes $\Delta n$ odd.
\end{proof}

Thus the squarefree part is $D=q$, and $\left(\frac{q}{p}\right)=1$.

\subsection{\texorpdfstring{The case $A=1$}{The case A=1}}

\begin{lemma}\label{lem:Aone}
If $A=1$, the two partitions are
\[
11=2^3+3=2+3^2.
\]
\end{lemma}
\begin{proof}
Write $\Delta m=2k$ with $k\geq1$. Then
\[
q^{n_1}=2^{2k}-1=(2^k-1)(2^k+1).
\]
The two factors are odd and differ by $2$, so they are coprime.
As their product is a power of a single prime, the smaller factor
must be $1$. Hence $k=1$, $q=3$, and $n_1=1$.
The other defining equation is now
\[
2^{m_2}=3^{\Delta n}-1.
\]
Since $\Delta n$ is odd, $3^{\Delta n}-1\equiv2\pmod8$.
Thus $m_2=1$, which gives $\Delta n=1$.
Therefore $m_1=3$, $n_2=2$, and $p=11$.
\end{proof}

\section{Valuations, size bounds, and orders}
\label{sec:bounds}

Assume $A>1$ from now on, and write
\[
M=\Delta m=2k,\qquad N=\Delta n\quad\text{(odd)}.
\]
All logarithms without a subscript are natural logarithms.
% Source: even_odd_case.tex, arithmetic of A and exponent bounds.

\subsection{\texorpdfstring{Exact valuations and the factorization of $A$}{Exact valuations and the factorization of A}}

\begin{lemma}\label{lem:valuations}
We have
\[
n_1=v_q(2^M-1),\qquad m_2=v_2(q-1).
\]
Writing
\[
q-1=2^{m_2}u,\qquad
\Psi=\frac{q^N-1}{q-1}=1+q+\cdots+q^{N-1},
\]
we have
\[
\boxed{A=u\Psi,\qquad 2^M-1=q^{n_1}u\Psi.}
\]
Here $u$ and $\Psi$ are odd, and $q$ divides neither of them.
\end{lemma}
\begin{proof}
The first valuation follows from $q\nmid A$.
Since $N$ is odd, $\Psi$ is a sum of an odd number of odd terms,
so $v_2(q^N-1)=v_2(q-1)$.
Apply this to $A2^{m_2}=q^N-1$, using that $A$ is odd.
The factorizations follow by substitution.
\end{proof}

Both $u$ and $\Psi$ divide $2^M-1$. They need not be coprime:
since $\Psi\equiv N\pmod{q-1}$,
\[
\gcd(u,\Psi)=\gcd(u,N).
\]
\subsection{\texorpdfstring{Size bounds and the exclusion of $N=1$}{Size bounds and the exclusion of N=1}}

\begin{lemma}\label{lem:bounds}
Write $M=2k$. Then
\[
q^{n_1}\leq2^k+1,\qquad A\geq2^k-1,\qquad m_2\leq k,
\]
and
\[
\boxed{2^{M/2}<q^N<2^M.}
\]
Moreover $N\geq3$, $q<2^{M/3}$, and $M\geq6$.
\end{lemma}
\begin{proof}
The factors in
\[
2^M-1=(2^k-1)(2^k+1)
\]
are coprime. Thus $q^{n_1}$ divides one factor, and the other divides
$A$. This proves the first two bounds. Also
$2^{m_2}\leq q-1\leq2^k$, proving the third.

For the upper bound, put $c=2^{m_2}$. Then
\[
q^N-1=\frac{c}{q^{n_1}}(2^M-1)
 \leq\frac{q-1}{q}(2^M-1)<2^M-1.
\]
For the lower bound,
\[
q^N=cA+1\geq2(2^k-1)+1>2^k.
\]

If $N=1$, then $A=u<q$. Hence
\[
2^{2k}-1=Aq^{n_1}<q(2^k+1),
\]
so $q>2^k-1$. Since $q$ is odd, $q\geq2^k+1$.
Together with $q\leq q^{n_1}\leq2^k+1$, this forces
$q=2^k+1$ and $n_1=1$. Then $q-1=2^k$ has odd part $u=1$,
so $A=1$, a contradiction.
Thus $N\geq3$, and $q^3\leq q^N<2^M$.
Since $q\geq3$ and $M$ is even, $M\geq6$.
\end{proof}

Equivalently,
\begin{equation}\label{eq:window}
\frac{M\log2}{2\log q}<N<\frac{M\log2}{\log q}.
\end{equation}
The interval has length proportional to $M/\log q$; these bounds alone
do not make it empty for large $M$.

The larger exponents satisfy
\[
m_1=\lfloor\log_2p\rfloor,\qquad
n_2=\lfloor\log_qp\rfloor.
\]
Indeed, $q^{n_1}\leq2^k+1<2^{m_1}$ gives
$2^{m_1}<p<2^{m_1+1}$. Also $q^{n_2}<p$, and, since
$c=2^{m_2}\leq q-1$ and $p>q$,
\[
q^{n_2+1}=q(p-c)>p.
\]

\subsection{Orders and multiplicities of the other prime factors}

\begin{lemma}\label{lem:orders}
For every prime $\ell\mid2^M-1$ with $\ell\neq q$,
\[
\operatorname{ord}_\ell(q)\mid N.
\]
Consequently the integer
\[
L=\operatorname{lcm}_{\substack{\ell\mid2^M-1\\
\ell\text{ prime},\ \ell\neq q}}
\operatorname{ord}_\ell(q)
\]
is odd and divides $N$. The lcm of an empty set is taken to be $1$.
\end{lemma}
\begin{proof}
Every such $\ell$ divides $A$, and $A\mid q^N-1$.
\end{proof}

Primes $\ell\mid q-1$ have order $1$ and can be omitted from the lcm.
Since $3\mid2^M-1$, it follows that
\[
q=3\quad\text{or}\quad q\equiv1\pmod3.
\]
Thus
\[
L\mid N,\qquad N\text{ odd},\qquad
2^{M/2}<q^N<2^M.
\]
For each odd prime $\ell\mid q-1$, put
$d_\ell=\operatorname{ord}_\ell(2)$. Since $\ell\mid u\mid2^M-1$,
$d_\ell\mid M$. Lifting the exponent on the two sides of
$2^M-1=Aq^{n_1}$ gives
\begin{equation}\label{eq:valuation-coupling}
\boxed{
v_\ell(2^{d_\ell}-1)+v_\ell(M)
 =v_\ell(q-1)+v_\ell(N).}
\end{equation}
Here $\ell\nmid d_\ell$ because $d_\ell\mid\ell-1$, and $\ell\neq q$.
For example, if $q\neq3$, then $3\mid q-1$ and
\[
1+v_3(M)=v_3(q-1)+v_3(N).
\]
Thus the order conditions and the multiplicity conditions constrain
the same two gaps.

\section{\texorpdfstring{Eliminating $q=3$ and restricting $q>3$}{Eliminating q=3 and restricting q>3}}
\label{sec:parities}
% Added in discussion: q=3 proof, strengthened parity reduction, and its
% consequences. These replace the provisional table in even_odd_case.tex.

\subsection{\texorpdfstring{The base $q=3$}{The base q=3}}

\begin{lemma}\label{lem:qthree}
For $q=3$, the only prime with two partitions is $p=11$.
\end{lemma}
\begin{proof}
The case $A=1$ is settled, so suppose $A>1$.
We have $m_2=1$, $N$ odd, and
\[
n_1=v_3(2^M-1)=1+v_3(M).
\]
If $n_1>1$, then $3\mid M$. Since $\operatorname{ord}_7(2)=3$,
we get $7\mid2^M-1$, hence $7\mid A$, hence $7\mid3^N-1$.
But $\operatorname{ord}_7(3)=6$, contradicting $N$ odd.
Thus $n_1=1$, and the original partitions give
\[
3^{n_2}-1=2^{m_1},\qquad n_2=1+N\text{ even}.
\]
The factors $3^{n_2/2}-1$ and $3^{n_2/2}+1$ are powers of $2$
differing by $2$, so they are $2$ and $4$. Thus $n_2=2$, $m_1=3$,
and $p=11$, which has $A=1$.
\end{proof}

\subsection{\texorpdfstring{Parities for $q>3$}{Parities for q>3}}

\begin{lemma}\label{lem:new-parities}
In the remaining case $A>1$, we have $q>3$ and
\[
m_1,m_2,n_1\text{ even},\qquad n_2\text{ odd}.
\]
Consequently $n_1\geq2$, $n_2\geq5$, and $q^2\mid2^M-1$.
\end{lemma}
\begin{proof}
For $q>3$ we know $q\equiv1\pmod3$. Then
$p=2^{m_i}+q^{n_i}\not\equiv0\pmod3$ forces both $m_i$ even.
Suppose $n_2$ is even. Since $m_1$ is even, the left-hand side of
\[
\bigl(q^{n_2/2}\bigr)^2-\bigl(2^{m_1/2}\bigr)^2
=q^{n_1}-2^{m_2}>0
\]
is a difference of distinct integer squares, and hence at least
$2\cdot2^{m_1/2}+1$. By Lemma~\ref{lem:bounds},
$q^{n_1}\leq2^k+1=2^{M/2}+1$. Hence the right-hand side is at most
$2^{M/2}+1-2^{m_2}$, which is smaller. Therefore $n_2$ is odd,
and $N$ odd makes $n_1$ even. Finally $n_2=n_1+N\geq2+3$.
\end{proof}

Both $q\equiv3$ and $q\equiv7\pmod8$ would give $m_2=1$, so both
are excluded. The remaining possibilities are
\[
\begin{array}{c|c}
q\bmod8&m_2=v_2(q-1)\\ \hline
5&2\\
1&\text{even and at least }4
\end{array}
\]
with $n_1$ even and $n_2$ odd in both rows.
Lemma~\ref{lem:q-one-mod-eight} below excludes the $5\pmod8$ row.
Also $q<p^{1/5}$ since $n_2\geq5$.
The size estimate also sharpens to
\begin{equation}\label{eq:sharp-A}
A<\frac{\sqrt{2^M q^N}}{2^{m_2/2}q^{n_1/2}}
\leq\frac{\sqrt{2^M q^N}}{2q},\qquad
A<\frac{p}{2^{m_2}q^{n_1}}\leq\frac{p}{4q^2}.
\end{equation}
The latter follows from $p=A2^{m_2}q^{n_1}+2^{m_2}+q^{n_1}$.


\section{Restrictions from Mersenne factors}
\label{sec:mersenne}

\subsection{\texorpdfstring{The order and multiplicity of $q$}{The order and multiplicity of q}}
% Source: original valuation discussion; the even multiplicity consequence
% and the development target use the added parity and q=3 reductions.
Put
\[
r=\operatorname{ord}_q(2),\qquad \rho=v_q(2^r-1).
\]
Then $q$ is a primitive Mersenne factor at exponent $r$: it divides
$2^r-1$ and no $2^a-1$ with $1\leq a<r$.
This describes its order, not its multiplicity. Since $r\mid M$ and
$q\nmid r$, lifting the exponent gives
\begin{equation}\label{eq:exact-q-valuation}
\boxed{n_1=\rho+v_q(M).}
\end{equation}
Thus if $q\nmid M$, then $\rho=n_1$ is even and at least $2$.
If $\rho=1$, then $v_q(M)$ is odd and positive, so $q\mid M$.
In either case $2^M\equiv1\pmod{q^2}$; if $q\nmid M$, also
$2^r\equiv2^{q-1}\equiv1\pmod{q^2}$.

The $q=3$ proof suggests a development target. A larger multiplicity of
$3$ forces $3\mid M$, which forces the Mersenne factor $7$; its order at
base $3$ then contradicts $N$ odd. For general $q$, seek a forced factor
$\ell\neq q$ of $2^M-1$ for which $\operatorname{ord}_\ell(q)$ cannot
divide $N$. More generally, several forced orders may have an lcm that
cannot fit the size interval \eqref{eq:window}. A uniform way to force
these incompatible orders is still missing.

\subsection{Orders modulo A and coprime gaps}
\label{sec:exact-orders}
% Added in discussion; no coprimality assumption in the first lemma.
\begin{lemma}
For $A>1$, $q>3$,
\[
\operatorname{ord}_A(2)=M,\qquad \operatorname{ord}_A(q)=N.
\]
\end{lemma}
\begin{proof}
Lemma~\ref{lem:bounds} gives $A\geq2^{M/2}-1$.
Equality would give $q^{n_1}=2^{M/2}+1$. Since $n_1$ is even,
$(q^{n_1/2}-1)(q^{n_1/2}+1)=2^{M/2}$ forces $q=3$.
Thus $A>2^{M/2}-1$. A proper divisor of $M$ is at most $M/2$,
so cannot be the order of $2$ modulo $A$.
Also
\[
A=\frac{q^N-1}{2^{m_2}}\geq\frac{q^N-1}{q-1}>q^{N-1},
\]
which proves the second assertion.
\end{proof}

Suppose $\gcd(M,N)=1$. The subgroups generated by $2$ and $q$
modulo $A$ have trivial intersection. Consequently
\[
\operatorname{ord}_A(2q)=MN,\qquad MN\mid\varphi(A).
\]
For each prime $\ell\mid A$, put
$r_\ell=\operatorname{ord}_\ell(2)$, $t_\ell=\operatorname{ord}_\ell(q)$.
Then $\gcd(r_\ell,t_\ell)=1$ and $r_\ell t_\ell\mid\ell-1$.
In particular, a primitive Mersenne factor $\ell\neq q$ at exponent $M$
has the form $\ell=hM+1$ and satisfies
\[
\boxed{\operatorname{ord}_\ell(q)\mid\gcd(h,N).}
\]
If $\gcd(h,N)=1$, then $\ell\mid q-1$. This excludes $\ell>q$,
in particular when $h$ is a power of $2$.
When $q\mid M$, every primitive factor at $M$ exceeds $q$;
its multiplier $h$ must therefore share an odd prime factor with $N$.

For positive integers $a,b$, the same subgroup argument prime power
by prime power gives
\[
\gcd(A,2^a q^b-1)=\gcd(A,2^a-1,q^b-1)
\quad\text{if }\gcd(M,N)=1.
\]

\subsection{Common factors of the two gaps}
% Source: diophantine_analysis.tex (divisor lemma) and even_odd_case.tex
% (common-gap construction). The q>3 specialization uses the added parities.

\begin{lemma}\label{lem:divisor}
If $d>1$ divides both gaps, write
\[
M=ad,\qquad N=bd,\qquad Y=2^aq^b.
\]
Then
\[
Ap=(Y-1)(1+Y+\cdots+Y^{d-1}),\qquad Y-1\mid A.
\]
\end{lemma}
\begin{proof}
The factorization follows from $Ap=Y^d-1$. If $p\mid Y-1$, then
\[
A=\frac{Y-1}{p}(1+Y+\cdots+Y^{d-1})
 \geq 1+Y+\cdots+Y^{d-1}>Y^{d-1}.
\]
But \eqref{eq:A-size} gives $A<Y^{d/2}/\sqrt6$, a contradiction
for $d\geq2$. Hence $p\nmid Y-1$, and $Y-1\mid Ap$ implies
$Y-1\mid A$.
\end{proof}

Whether $\gcd(M,N)=1$ remains open. This is a separate question from
whether $q\mid M$.
Let an odd prime $s$ divide both gaps, and write
\[
U=2^{M/s},\quad V=q^{N/s},\quad Y=UV,\quad B=Y-1.
\]
Then
\[
Ap=(Y-1)\Phi_s(Y),\qquad
\Phi_s(Y)=1+Y+\cdots+Y^{s-1}.
\]
Lemma~\ref{lem:divisor} gives $p\nmid B$ and $B\mid A$, so
\[
B\mid2^M-1,\quad B\mid q^N-1,\quad
p\mid\Phi_s(Y),\quad \operatorname{ord}_p(Y)=s,\quad p\equiv1\pmod s.
\]
Thus a common factor of the gaps forces a particular divisor $B$ into
$A$, not merely an unspecified factor of a Mersenne number.

The case $s=3$ is impossible. Indeed $B\mid U^3-1,V^3-1$.
With $W=\min(U,V)$ and $Z=\max(U,V)$, this gives
$W^2\equiv Z\pmod{WZ-1}$. Both are positive and smaller than $WZ-1$,
so $W^2=Z$, impossible for powers of distinct primes. Hence every prime
factor of $\gcd(M,N)>1$ is at least $5$.

For $s\geq5$, reduce $V^s-1$ modulo $B=UV-1$ to obtain
$B\mid V^{s-1}-U$. Also $V>\sqrt U$ by $q^N>2^{M/2}$,
so $V^{s-1}>U$. Thus there is an integer $\tau\geq1$ with
\[
V^{s-1}-U=\tau(UV-1),\qquad
U(\tau V+1)=V^{s-1}+\tau.
\]
Here $U$ is a square, and $\tau$ is odd by parity. Since the remaining
$q$ satisfy $q\equiv1\pmod4$ and $N/s$ is odd, put
$w=(V+1)/2$, which is odd. Reduction modulo $V+1$ gives
\[
U\frac{\tau-1}{2}\equiv-\frac{\tau+1}{2}\pmod w.
\]
The two consecutive integers $(\tau-1)/2$ and $(\tau+1)/2$ are coprime;
the congruence makes the first coprime to $w$. Consequently
\[
U\equiv-\frac{\tau+1}{\tau-1}\pmod w,
\qquad
\left(\frac{1-\tau^2}{\ell}\right)=1
\quad\text{for every prime }\ell\mid w.
\]
The last assertion uses that $U$ is a unit square modulo $w$.
In particular it holds for every prime dividing $(q+1)/2$.
These necessary conditions have not yet supplied a contradiction.

\subsection{Forbidden periods and forced residues}
\label{sec:forbidden-periods}
% Added in discussion: use minimum return periods and the exclusion of 3
% from N, without assuming that the gaps are coprime.
The exclusion of $3\mid\gcd(M,N)$ and the valuation identity at $3$
give
\[
3\nmid N,\qquad v_3(M)=v_3(q-1)-1.
\]
Indeed, if $3\mid N$, then $3\nmid M$, whereas
$1=v_3(q-1)+v_3(N)\ge2$.
Since $r=\operatorname{ord}_q(2)$ divides both $M$ and $q-1$,
\begin{equation}\label{eq:cubic-return}
\boxed{r\mid(q-1)/3,\qquad 2^{(q-1)/3}\equiv1\pmod q.}
\end{equation}
This is a necessary condition on $q$ alone.

If $\ell\mid A$ and $\ell-1=2^i3^j$, then
$\operatorname{ord}_\ell(q)\mid\gcd(N,\ell-1)=1$. Hence
\[
q\equiv1\pmod\ell.
\]
In particular $3\mid M$ forces $q\equiv1\pmod7$;
$12\mid M$, $q\ne13$, forces $q\equiv1\pmod{13}$;
and $18\mid M$ forces $q\equiv1\pmod{19}$.
The bases $7,19$ themselves are already excluded by $q\equiv1\pmod4$.
For example, $q=109$ passes \eqref{eq:cubic-return}, but
$\operatorname{ord}_{109}(2)=36$ forces $3\mid M$, while
$109\equiv4\pmod7$. Thus it is excluded for all $p$.

More generally, let $L\mid M$ be forced, and let $\mathcal S$ be
any set of primes $\ell\ne q$ such that
\[
\operatorname{ord}_\ell(2)\mid L,\qquad \ell-1=2^i3^j.
\]
Then, writing $P=\prod_{\ell\in\mathcal S}\ell$ and $c=m_2$,
\begin{equation}\label{eq:forced-residue-product}
\boxed{P\mid q-1,\qquad 2^cP\le q-1.}
\end{equation}
A single failed residue excludes $q$; the product can exclude a whole
specified family without computing the individual residues.

When $c=2$, $q\equiv5\pmod8$ implies that $2$ is a quadratic
nonresidue modulo $q$, so $4\mid r\mid M$. Thus $q\equiv1\pmod5$,
and the earlier classes modulo $48$ sharpen to
\[
q\equiv61\text{ or }181\pmod{240}\qquad(c=2).
\]

\subsection{The full common divisor}
\label{sec:full-gcd}
% Restored from even_odd_case.tex, Section 6.3; totient restriction added here.
Let $d=\gcd(M,N)>1$, $U=2^{M/d}$, $V=q^{N/d}$, and $Y=UV$.
Then $Ap=Y^d-1$. The sharper size bound gives
\[
p>2q(1-Y^{-d})Y^{d/2}>Y^{d/2}.
\]
For any proper divisor $e$ of $d$, $\Phi_e(Y)\mid Y^e-1<Y^{d/2}$.
Hence $p$ divides only the factor $\Phi_d(Y)$, and
\[
H_d:=\frac{Y^d-1}{\Phi_d(Y)}\mid A,\qquad
\operatorname{ord}_p(Y)=d,\qquad p\equiv1\pmod d.
\]
The M\"obius product gives
\[
\left|\log\Phi_d(Y)-\varphi(d)\log Y\right|
\leq\sum_{j\geq1}-\log(1-Y^{-j})
\leq\frac{Y}{(Y-1)^2}<\log2.
\]
Thus $\Phi_d(Y)<2Y^{\varphi(d)}$. Comparing the resulting lower bound
on $H_d$ with $A<Y^{d/2}/(2q)$ proves
\[
\boxed{\varphi(d)>d/2.}
\]
Equivalently, the prime factors of $d$ must satisfy
$\prod_{\ell\mid d}(1-1/\ell)>1/2$; every odd prime power passes this test.
Also $3\nmid d$ by the preceding argument.

\subsection{The proposed descent}
\label{sec:descent}
% User's proposed direction, worked out here. The obstruction below is to
% the reduced integrality system, not a counterexample to the original problem.
Here let $d\geq5$ be a prime dividing both gaps, and set
\[
a=M/d,\quad b=N/d,\quad U=2^a,\quad V=q^b,\quad B=UV-1.
\]
The gap bounds give $\sqrt U<V<U$ and $B\mid A$.
There are integers $\tau,\sigma\geq1$ such that
\begin{equation}\label{eq:quotient-pair}
V^{d-1}-U=\tau B,\qquad U^{d-1}-V=\sigma B.
\end{equation}
Both quotients are odd. Neither is $1$: for example $\tau=1$ would
require $V+1\mid V^{d-1}+1$, hence $V+1\mid2$.
Thus $\tau,\sigma\geq3$.

For fixed $\tau$, integrality is equivalent to
\begin{equation}\label{eq:quotient-divisor}
U=\frac{V^{d-1}+\tau}{\tau V+1}\in\mathbb Z
\quad\Longleftrightarrow\quad
E_d:=\frac{\tau^d+1}{\tau V+1}\in\mathbb Z.
\end{equation}
Indeed $\gcd(\tau,\tau V+1)=1$, and multiplying the first numerator
by $\tau^{d-1}$ gives residue $1+\tau^d$ modulo $\tau V+1$.
Equivalently choose a divisor $D\mid\tau^d+1$ with
$D>\tau+1$ and $D\equiv1\pmod\tau$, and put $V=(D-1)/\tau$.
Since $d$ is odd,
\[
\tau^d+1=(\tau+1)\Phi_{2d}(\tau),\qquad
\gcd(\tau+1,\Phi_{2d}(\tau))\mid d.
\]
This parametrizes integer quotients for fixed $\tau$; $\tau$ is unbounded.
The integer $E_d$ is not in general the cofactor $C=A/B$.

For the base-$V$ expansion, write
$U=\sum_j u_jV^j$, $\tau=\sum_j t_jV^j$ with digits in $[0,V-1]$.
The equation $U+\tau VU=V^{d-1}+\tau$ gives the carry relations
\[
u_j+\tau u_{j-1}+c_j=t_j+\mathbf1_{j=d-1}+Vc_{j+1},
\qquad u_{-1}=c_0=0,
\]
with zero terminal carry. In particular $U\equiv\tau\pmod V$.
These relations still require a restriction on the carries to yield an exclusion.

Put $c=v_2(V-1)=m_2$. The power-of-two condition on $U$ requires
\[
v_2(V^{d-1}+\tau)-v_2(\tau V+1)=a,
\]
and equality of the odd parts of the two expressions.
If $a>c+v_2(d-1)$, reduction modulo $2^a$ and lifting give
\[
v_2(\tau+1)=c+v_2(d-1),\quad
v_2(\tau V+1)=c,\quad v_2(E_d)=v_2(d-1).
\]
Thus the $2$-adic valuation of $E_d$ is controlled, but its odd part need
not be $1$. The congruence modulo $V+1$ proved above remains necessary;
it does not require the odd factors of $V+1$ to divide $U$.
The symmetric reduction modulo $U+1$ gives
\[
\gcd\left(\frac{\sigma-1}{2},U+1\right)=1,\qquad
V\equiv-\frac{\sigma+1}{\sigma-1}\pmod{U+1}.
\]
The gcd assertion follows from the coprimality of
$(\sigma-1)/2$ and $(\sigma+1)/2$.

The reduced conditions admit prime-power pairs. For example
\[
d=5,\quad U=4096=2^{12},\quad V=461\text{ prime},\quad
B=1888255,
\]
\[
\tau=23919,\qquad\sigma=149066189,\qquad
E_5=710020863679660.
\]
They satisfy \eqref{eq:quotient-pair}, $\sqrt U<V<U$, and
$v_2(E_5)=2$. Here $\tau>U$, so replacing $U$ by $\tau$ does not descend.
This pair fails the original problem: $461\equiv2\pmod3$ and
$2^{60}\equiv196\pmod{461}$.

A descent proving coprimality must therefore preserve more than
\eqref{eq:quotient-pair}. In particular it must retain
\[
q\mid U^d-1,\qquad n_1=v_q(U^d-1)\text{ even},\qquad
\frac{U^d-1}{q^{n_1}}=\frac{V^d-1}{2^{v_2(q-1)}},
\]
as well as the prime quotient $p=(U^dV^d-1)/A$.
No decreasing map preserving these conditions has yet been established.

\subsection{\texorpdfstring{The family $M=2s$, with $s$ prime}{The family M=2s, with s prime}}
% Added in discussion: this family has coprime gaps and q does not divide M.

Suppose $M=2s$ with $s$ an odd prime. Then $r\mid2s$;
since $q>3$, $r\neq1,2$, so $r=s$ or $2s$.
Thus $q\geq2s+1$ and
\[
N<\frac{2s\log2}{\log q}<s.
\]
Therefore $\gcd(M,N)=1$ and $q\nmid M$ in this family.
It belongs to the repeated primitive-factor case $\rho=n_1\geq2$ even;
the common-divisor construction above does not apply to it.

\section{The logarithmic bound}
\label{sec:logarithms}
% Added in discussion: the upper bound for Lambda, application of Laurent,
% and finite reduction for q dividing M. Laurent's estimate is external.

\subsection{Closeness of the two large powers}

Define
\[
\Lambda=n_2\log q-m_1\log2.
\]
The subtracted equation yields
\[
2^{m_1}(1-2^{-M})=q^{n_2}(1-q^{-N}),
\]
hence
\[
\Lambda=\log(1-2^{-M})-\log(1-q^{-N}).
\]
Since $q^N<2^M$ and $q^N\geq27$,
\begin{equation}\label{eq:lambda-upper}
\boxed{0<\Lambda<\frac{27}{26}q^{-N}.}
\end{equation}
Indeed, $-\log(1-x)\leq x/(1-x)$ for $0<x<1$.
In particular the linear form is nonzero.

Since $3\mid q-1$, its odd part $u$ is at least $3$. The gap equation
and $n_1\geq2$ give
\begin{equation}\label{eq:gap-ratio}
q^N=\frac{2^{m_2}}{q^{n_1}}(2^M-1)+1
 <\frac{2^M-1}{3q}+1,
\qquad \frac{2^{m_2}}{q^{n_1}}\leq\frac{q-1}{3q^2}.
\end{equation}
Along any hypothetical sequence with $q\to\infty$, this gives
$q^N/2^M\to0$. Meanwhile
\[
0<\frac{p}{2^{m_1}}-1
 \leq\frac{2^{M/2}+1}{2^{M+m_2}}\longrightarrow0
\]
as $M\to\infty$. The large terms $2^{m_1}$ and $q^{n_2}$ approach
each other; the gap powers need not.

\subsection{\texorpdfstring{An effective bound for fixed $q$}{An effective bound for fixed q}}

Miyazaki--Scott--Styer \cite[Proposition~6.1]{MSS}
(\texttt{prop-comp} in the local source) prompted the use of
Laurent's estimate here. Laurent \cite[Corollary~2]{Laurent} gives
\[
\log|b_2\log\alpha_2-b_1\log\alpha_1|
>-25.2H_1H_2\bigl(\max\{\log b'+0.38,10\}\bigr)^2,
\qquad b'=b_1/H_2+b_2/H_1,
\]
under its height and independence hypotheses.
Use $(\alpha_1,\alpha_2)=(2,q)$, $(b_1,b_2)=(m_1,n_2)$.
The bases are multiplicatively independent because they are distinct primes.
For a conservative application choose $H_1=H_2=\log q>1$;
these dominate both rational heights, both logarithms, and $1$.
This choice also avoids any issue with the mismatched index in the height
condition printed in the local copy. The asymmetric heights in Lemma~\ref{lem:uniform-N} give an absolute
bound on $N$. We keep this weaker estimate for the earlier finite check.

Put $\lambda=\log q$ and
\[
\mathcal H=\max\left\{\log\frac{m_1+n_2}{\lambda}+0.38,10\right\}.
\]
Combining the quoted lower bound with \eqref{eq:lambda-upper} gives
\begin{equation}\label{eq:baker-gap}
N\lambda<25.2\lambda^2\mathcal H^2+\log(27/26).
\end{equation}

The dependence on $m_1,n_2$ can be removed in favor of $M,q$.
Since $m_2\leq M/2$, $m_1\leq3M/2$. Also
$q^{n_2}<p<2^{m_1+1}$, so
\[
n_2<\frac{(3M/2+1)\log2}{\lambda}.
\]
Since $\lambda\geq\log3$, $\log2/\log3<2/3$, and $M\geq6$,
these inequalities give $m_1+n_2<3M$.
Together with $N\lambda>M\log2/2$, this proves
\begin{equation}\label{eq:baker-M}
\boxed{
\frac{M\log2}{2}
<25.2(\log q)^2
\left(\max\left\{\log\frac{3M}{\log q}+0.38,10\right\}\right)^2
+\log(27/26).}
\end{equation}

For each fixed $q$, the right side grows like $(\log M)^2$ whereas the
left side grows like $M$. It therefore gives an effective upper bound
on $M$ for that fixed $q$, and then bounds all four exponents by the
identities and bounds in Section~\ref{sec:bounds}. By itself this does not bound $q$
absolutely; the next step supplies such a bound in one branch.

\subsection{Minimum returns against the height bound}
\label{sec:minimum-returns}
% Added from the user's minimum-size/periodicity direction. These are
% necessary conditions, not a density assumption about prime powers.
Fix $q$ and a prime $d\ge5$ required to divide both gaps. Put
\[
c=v_2(q-1),\quad r=\operatorname{ord}_q(2),\quad
\rho=v_q(2^r-1),\quad s=v_3(q-1)-1.
\]
First require $c$ even and $r\mid(q-1)/3$.
Set $D=\operatorname{lcm}(2,d,r,3^s)$, and let $h$ be the least
integer with
\[
h\ge v_q(D),\qquad \rho+h\text{ even}.
\]
Define $L=Dq^{h-v_q(D)}$ and $e_0=\rho+h$. Then
\begin{equation}\label{eq:minimum-return}
\boxed{L\mid M,\qquad n_1\ge e_0\ge2,\qquad d\mid N.}
\end{equation}
This follows from $n_1=\rho+v_q(M)$ and
$v_3(M)=s$. More precisely $M=Lj$ requires
$3\nmid j$ and $v_q(j)$ even.
For $\rho=1$, $L=\operatorname{lcm}(2,d,qr,3^s)$.
Thus even the minimum permitted $n_1$ requires a factor $q$ in $M$.
If $d=q$ and $\rho$ is even, the minimum is instead $v_q(M)=2$.

The minimum sizes give
\begin{equation}\label{eq:return-size}
2^M\ge1+\frac{q^{e_0}(q^d-1)}{2^c}.
\end{equation}
If $T$ is the base-$2$ logarithm of the right side, then
$M\ge L\lceil T/L\rceil$. This lower bound must fit below the
effective upper bound from \eqref{eq:baker-M}.
The inequalities alone do not exclude all $q$ when $\rho$ is even.

\subsubsection*{Propagating the required periods}
Start with forced divisors $L\mid M$, $K\mid N$, initially the $L$
above and $K=d$. For each selected prime power $\ell^f\mid2^L-1$,
$\ell\ne q$, the full gap equations force $\ell^f\mid A$. Therefore
replace $K$ by
\[
\operatorname{lcm}\bigl(K,\operatorname{ord}_{\ell^f}(q)\bigr).
\]
If $\gcd(K,6)>1$, there is a contradiction.
Conversely, for each odd prime $\ell\mid q-1$, put
$r_\ell=\operatorname{ord}_\ell(2)$ and
$\eta_\ell=v_\ell(2^{r_\ell}-1)$. The valuation identity forces
\[
r_\ell\mid M,\qquad
v_\ell(M)\ge v_\ell(q-1)+v_\ell(K)-\eta_\ell.
\]
Enlarge $L$ by these forced factors, retaining the exact $3$-valuation
and the parity condition on $v_q(M)$. Newly forced Mersenne factors
can then enlarge $K$ again. All these implications hold for every
multiple of the current $L,K$.
Use the resulting lower bound on $n_1$ and replace $d$ by $K$ in
\eqref{eq:return-size}. A contradiction follows if this forces $M$
beyond \eqref{eq:baker-M}, or if \eqref{eq:forced-residue-product} fails.
The open step is to force such a contradiction uniformly, rather
than merely obtain compatible divisibility conditions.

The doubling-map formulation in \cite[Theorem A and \S4]{AGJ26}
identifies the period of $1/Q$ under doubling with
$\operatorname{ord}_Q(2)$, for odd $Q$. At prime-power denominators,
lifting gives
\[
\operatorname{ord}_{q^e}(2)=r q^{\max(0,e-\rho)}.
\]
This expresses the minimum-return obstruction dynamically. The paper's
compressed returns compute these periods; an additional restriction on
the simultaneous periods at bases $2$ and $q$ is still required here.

\subsection{\texorpdfstring{An effective reduction when $q\mid M$}{An effective reduction when q divides  M}}

If $q\mid M$, then $qr\mid M$ because $r\mid M$ and $q\nmid r$.
Since $q<2^r$, this gives
\[
M\geq qr>\frac{q\log q}{\log2}.
\]
Together with \eqref{eq:baker-M}, this bounds $q$ absolutely in this branch.
A coarse cutoff is $q<10^6$.

Put $\lambda=\log q$, $t=M/\lambda$, and
$H(t)=\max\{\log(3t)+0.38,10\}$. Dividing \eqref{eq:baker-M} gives
\[
\frac{t}{H(t)^2}
 <\frac{50.4}{\log2}\lambda
   +\frac{2\log(27/26)}{\lambda\log2\,H(t)^2}
 <73\lambda.
\]
If $q\geq10^6$, then $t>q/\log2$, $H(t)=\log(3t)+0.38>2$,
and $t/H(t)^2$ is increasing. Hence
\[
q<73\log2\,\lambda(\lambda+C)^2,
\qquad C=\log(3/\log2)+0.38<1.845126.
\]
At $q=10^6$ the right side is less than $172000$.
Moreover $q/[\log q(\log q+C)^2]$ is increasing thereafter:
its logarithmic derivative with respect to $\lambda$ is
$1-1/\lambda-2/(\lambda+C)>0$.
Thus $q<10^6$ in this branch. For each such $q$,
\eqref{eq:baker-M} bounds $M$ effectively. The uniform bound and
finite exclusion below complete this branch.

\begin{lemma}
If $q\mid M$, then $36\nmid M$ and $48\nmid M$.
\end{lemma}
\begin{proof}
If $36\mid M$, the primes $3,5,7,13,19,37$ divide $2^M-1$;
their orders at base $2$ are $2,4,3,12,18,36$ respectively.
Each has $\ell-1$ supported on $2,3$, so
Section~\ref{sec:forbidden-periods} forces $q\equiv1\pmod\ell$
unless $q=\ell$. The bases $3,5,7,19$ are already excluded.
Since $4\mid M$, $q\equiv1\pmod5$, excluding $q=13,37$ as well.
Consequently
\[
959595=3\cdot5\cdot7\cdot13\cdot19\cdot37\mid q-1.
\]
As $4\mid q-1$, this gives $q\ge3838381$, contrary to $q<10^6$.

If $48\mid M$, use $3,5,7,13,17,97$, whose orders are
$2,4,3,12,8,48$. The base $17$ is excluded modulo $3$;
$13,97$ are excluded by $q\equiv1\pmod5$.
Thus $2250885\mid q-1$, giving $q\ge9003541$.
\end{proof}
In particular, in the $q\mid M$ branch,
\[
\operatorname{lcm}(2,r,3^{v_3(q-1)-1})
\quad\text{cannot be divisible by }36\text{ or }48.
\]
These exclusions apply to all $p$; they use no bound on $p$ and no
enumeration. Additional prime factors in the residue product can
strengthen the thresholds.

\subsection{Excluding the branch with a finite check}
\label{sec:q-divides-exclusion}
% New computer-assisted deduction from the minimum-return argument.
% This subsumes the earlier special-family check_power_two_orders.lua.
\begin{lemma}
If $q\mid M$, then $M<4\cdot10^6$.
\end{lemma}
\begin{proof}
Write $H_\lambda(x)=\max\{\log(3x/\lambda)+0.38,10\}$.
The preceding estimates give $M<73\lambda^2H_\lambda(M)^2$.
Since $\lambda=\log q<14$ and $\lambda H_\lambda(M)$ is increasing
in $\lambda$,
\[
\frac{M}{H_{14}(M)^2}<73\cdot14^2=14308.
\]
The function $x/H_{14}(x)^2$ is increasing: on the logarithmic
branch its derivative has the sign of $H_{14}(x)-2>0$.
But $H_{14}(4\cdot10^6)<15$, and
$4\cdot10^6/15^2>14308$. Hence $M<4\cdot10^6$.
\end{proof}

\begin{lemma}\label{lem:q-not-M}
For $A>1$ and $q>3$, $q\nmid M$.
\end{lemma}
\begin{proof}[Proof with finite computation]
Suppose $q\mid M$. It suffices to check primes $3<q<10^6$ satisfying
\[
q\equiv1\pmod3,\quad v_2(q-1)\text{ even},\quad
2^{(q-1)/3}\equiv1\pmod q.
\]
For each such prime, put $r=\operatorname{ord}_q(2)$,
$s=v_3(q-1)-1$, and
\[
L=\operatorname{lcm}\left(
q\operatorname{lcm}(2,r,3^s),\
\{\operatorname{ord}_\ell(2):\ell\mid q-1,\ \ell\text{ odd prime}\}
\right).
\]
Every entry in this lcm divides $M$: the odd part of $q-1$ divides
$A$. Thus reject $q$ if $L\ge4\cdot10^6$ or $v_3(L)>s$.
For the remaining primes, the following implications are necessary:
\[
4\mid L\Rightarrow q\equiv1\pmod5,\qquad
3\mid L\Rightarrow q\equiv1\pmod7,\qquad
5\mid L\Rightarrow q^5\equiv1\pmod{31}.
\]
They follow from $\gcd(N,6)=1$ and the respective orders
$\operatorname{ord}_5(2)=4$, $\operatorname{ord}_7(2)=3$,
$\operatorname{ord}_{31}(2)=5$. The bases $5,7,31$ fail the
initial restrictions; the script also explicitly omits a witness equal to $q$.

The exact sieve in \path{check_q_divides_gap.lua} gives:
\[
\begin{array}{l|r}
\text{primes passing the initial restrictions}&4347\\\hline
\text{excluded by }L\ge4\cdot10^6&4332\\
\text{excluded by }v_3(L)>s&5\\
\text{excluded modulo }5&7\\
\text{excluded modulo }7&2\\
\text{excluded modulo }31&1\\\hline
\text{remaining}&0
\end{array}
\]
These are successive, disjoint rejection counts. The script enumerates
only $q$ and uses integer arithmetic; it does not enumerate $p,M,N$.
Its arithmetic stays below $2^{53}$, including on LuaJIT.
\end{proof}

Consequently every remaining solution requires
\[
\boxed{q\nmid M,\qquad n_1=\rho\ge2\text{ even},\qquad
q^{n_1}\mid2^r-1,\qquad r\mid(q-1)/3.}
\]
In particular $2^{(q-1)/3}\equiv1\pmod{q^2}$.
A prime $d\mid\gcd(M,N)$ must be different from $q$.
This leaves the repeated primitive-factor branch; coprimality of
the two gaps remains open.

\section{The smaller-term ratio and common divisors}
\label{sec:additional-returns}
% New deductions following the user's minimum-return and common-gap split.
Throughout this section $q\nmid M$, by Lemma~\ref{lem:q-not-M}.
Put $e=n_1=\rho$, $c=m_2$, $Q=q^e$, and
\[
E=\operatorname{lcm}(2,r),\qquad M=hE.
\]
Then $Q\mid2^r-1$, and lifting gives
\[
A=\frac{2^r-1}{Q}\,\frac{2^M-1}{2^r-1},\qquad
\gcd\left(q,\frac{2^M-1}{2^r-1}\right)=1.
\]
In the cyclotomic factorization of $2^M-1$, only $\Phi_r(2)$ contains
$q$. Every other factor $\Phi_j(2)$, $j\mid M$, enters $A$ intact.

\subsection{Congruences on p}
\label{sec:p-common-gap}
Let $g=\gcd(M,N)>1$. Section~\ref{sec:full-gcd} gives
$p\equiv1\pmod g$. If a prime $d\mid g$ also divided $q-1$,
then $p=2^{m_2}+q^{n_2}$ would give
$1\equiv2^{m_2}+1\pmod d$, impossible. Thus
\[
\gcd(g,q-1)=1.
\]
Also $p=2^{m_1}+q^{n_1}$ is a sum of two nonzero squares modulo $5$.
Since $p>5$ and $q\ne5$, $p\equiv2$ or $3\pmod5$.
Consequently $5\nmid g$. Together with the previous exclusions,
\begin{equation}\label{eq:common-gap-coprimality}
\boxed{\gcd(g,30q(q-1))=1,\qquad \gcd(g,E)=1,\qquad gE\mid M.}
\end{equation}
Every prime factor of $g$ is at least $7$.
The case in which a common prime has its full multiplicity already
in $r$ is therefore excluded: such a prime would divide $q-1$.

More generally, let $\ell$ be an odd prime dividing $q-1$, and put
$\eta_\ell=v_\ell(2^{\operatorname{ord}_\ell(2)}-1)$.
If $\eta_\ell=1$, the valuation identity and
$\gcd(g,q-1)=1$ imply
\begin{equation}\label{eq:ordinary-q-minus-one}
\boxed{\ell\nmid N,\qquad v_\ell(M)=v_\ell(q-1)-1.}
\end{equation}
Indeed, $\ell\mid N$ would make
$v_\ell(M)=v_\ell(q-1)+v_\ell(N)-1\ge1$, contradicting
$\gcd(g,q-1)=1$. The same conclusion holds when $g=1$.
Since $r\mid M$, we obtain the stronger condition on $q$
\[
r\mid(q-1)/\ell,\qquad
2^{(q-1)/\ell}\equiv1\pmod{q^2}.
\]
This applies to $\ell=3,5,7,11,13$ whenever $\ell\mid q-1$.

\subsection{The ratio before taking d-th powers}
For any common divisor $d>1$, write $M=ad$, $N=bd$ and
$U=2^a$, $V=q^b$. Keep the smaller terms in the equality:
\[
p=2^{m_2}U^d+q^{n_1}=2^{m_2}+q^{n_1}V^d,
\qquad
R:=\frac{q^{n_1}}{2^{m_2}}=\frac{U^d-1}{V^d-1}.
\]
Thus $1<U/V<R^{1/d}$. More precisely, with
$\delta=a\log2-b\log q>0$,
\begin{equation}\label{eq:smaller-ratio}
\boxed{d\delta=\log R-\Lambda,\qquad
0<\Lambda<\frac{27}{26}q^{-bd}.}
\end{equation}
For fixed $q$, both $n_1$ and $m_2$ are fixed, so $R$ is fixed.
Increasing $d$ requires a closer ratio $2^a/q^b$ and the
exponentially accurate equality \eqref{eq:smaller-ratio}.

The congruences on $q$ also constrain $R$. With
$u=(q-1)/2^{m_2}$,
\[
R=\frac{u q^{n_1}}{q-1},\qquad
\log R=(n_1-1)\log q+\log u+\log\frac q{q-1},
\]
where $3\mid u$ and $0<\log(q/(q-1))<1/(q-1)$.
Consequently
\[
\frac{(n_1-1)\log q+\log u-(27/26)q^{-N}}d
<\delta<
\frac{(n_1-1)\log q+\log u+1/(q-1)}d.
\]
Every odd product forced to divide $q-1$ also divides $u$, giving
a corresponding lower bound on this ratio. For a common prime $d$,
the congruence on $p$ gives, in addition,
\[
2^{m_2+a}+q^{n_1}\equiv1\pmod d,\qquad
2^{m_2}+q^{n_1+b}\equiv1\pmod d.
\]
These follow by reducing the $d$-th powers modulo $d$.

\subsection{Height using the least even multiple of r}
The coprime factors $2^{E/2}-1$ and $2^{E/2}+1$ give
\[
Q\le2^{E/2}+1,\qquad c\le E/2,\qquad
E\ge2\log_2(q^2-1).
\]
Writing $x=2^{E/2}$, $y=2^{M-E/2}$, the gap equation gives
\[
q^N\ge\frac{2(2^M-1)}{2^{E/2}+1}+1
=\frac{2xy+x-1}{x+1}>y.
\]
Consequently
\begin{equation}\label{eq:additional-height}
\boxed{N\log q>\left(1-\frac1{2h}\right)M\log2.}
\end{equation}
Also $m_1\le M+E/2$ and
$n_2<(m_1+1)\log2/\log q$. Since $q\ge13$ and $M\ge6$,
these imply $m_1+n_2<2M$; if $h\ge5$, they give
$m_1+n_2<3M/2$.
Combining with \eqref{eq:baker-gap}, for $h\ge5$ we obtain
\begin{equation}\label{eq:baker-additional}
\left(1-\frac1{2h}\right)M\log2
<25.2(\log q)^2
\left(\max\left\{\log\frac{3M}{2\log q}+0.38,10\right\}\right)^2
+\log(27/26).
\end{equation}
For $h<5$, replace $3M/2$ inside the logarithm by $2M$.
These improve the upper bound for fixed $q$; they do not yet give
an absolute bound on $q$.

The multiplier $h=M/E$ also has an elementary upper bound in terms of
$N$. Since $2^{m_1}<q^{n_2}$ and
$E\log2\ge2\log(Q-1)>2(e-1/4)\log q$,
\begin{equation}\label{eq:number-of-returns}
\boxed{h<\frac{N+e}{2e-1/2}\le\frac{2(N+2)}7.}
\end{equation}
The last fraction is obtained at the minimum $e=2$; the middle
expression decreases with $e$. In particular $N=5$ forces $h=1$,
so $M=E$. For $N=7$ we have $h\le2$, and for $N=11$, $h\le3$.
Since $q^{n_2}<p<2q^{n_2}$, we have
$N+e=n_2=\lfloor\log_q p\rfloor$.
Lemma~\ref{lem:sharp-h} improves the middle bound to $h<n_2/(2e)$.

\subsection{The factor forced by a common prime}
Let $d\ge7$ be prime with $d\mid\gcd(M,N)$, and write
\[
M=ad,\quad N=bd,\quad U=2^a,\quad V=q^b,\quad
B=UV-1,\quad C=A/B.
\]
By \eqref{eq:common-gap-coprimality}, $E\mid a$.
Then $Q\mid U-1$, and $d\ne q$ gives
\begin{equation}\label{eq:intact-cyclotomic}
A=\frac{U-1}{Q}\Phi_d(U),\qquad
\gcd(q,\Phi_d(U))=1,\qquad \Phi_d(U)\mid A.
\end{equation}
Thus $h\ge d$, and the minimum sizes and height bound imply
\[
M\ge2d\log_2(q^2-1),\qquad
N\log q>\left(1-\frac1{2d}\right)M\log2.
\]
There is also a restriction on the other gap:
\begin{equation}\label{eq:five-d}
\boxed{N\ge5d.}
\end{equation}
Indeed $U>Q=q^e$, whereas
$U^d=1+Q(V^d-1)/2^c<QV^d$.
Hence $b>(1-1/d)e\ge8/5$.
Since $b$ is odd and $3\nmid b$, it follows that $b\ge5$.
Using that $U$ is a square gives the stronger inequality
\[
b>\left(2-\frac1d\right)e-\frac12:
\]
$Q\le\sqrt U+1$ and $q^e-1>q^{e-1/4}$ imply
$U>q^{2e-1/2}$, which can be substituted into $U^d<QV^d$.
In particular, $N=5d$ forces $n_1=e=2$.
Since $d\mid h$ and \eqref{eq:number-of-returns} gives
$h<(10d+4)/7<2d$, it also forces $h=d$. Thus
\[
N=5d\quad\Longrightarrow\quad n_1=2,\qquad M=dE.
\]

The complete factor in \eqref{eq:intact-cyclotomic} supplies
\[
\frac{\Phi_d(U)}{\gcd(\Phi_d(U),UV-1)}\mid C,
\qquad
C<\frac{(UV)^{d/2}}{2^{c/2}q^{e/2}(UV-1)}.
\]
For every prime power $\ell^f\mid\Phi_d(U)$,
$\operatorname{ord}_{\ell^f}(q)\mid N$.
This is the full factor to which the period and residue restrictions
can be applied; the mixed-base gcd still needs control.

\subsection{Consequences for the full common divisor}
The arguments giving $N\ge5d$ and $N=5d\Rightarrow n_1=2$, $M=dE$
use only $dE\mid M$, $d\mid N$, $d\ne q$, and the gap equations.
They therefore apply with $d=g=\gcd(M,N)$ as well, even when $g$
is composite. Thus
\[
g>1\quad\Longrightarrow\quad N\ge5g\ge35.
\]
In particular the gaps are coprime when $N$ is prime, or when $N$
is a power of $5$.

The first possible value $N=35$ is also excluded from the noncoprime
case. It would require $g=7$, $n_1=2$ and $n_2=37$.
Since $p\equiv1\pmod7$, its partition as a sum of two nonzero
squares gives
\[
2^{m_1}\equiv q^2\equiv4\pmod7.
\]
Hence $q\equiv2$ or $5\pmod7$. The other partition gives
$2^{m_2}\equiv1-q^{37}\equiv1-q\pmod7$, namely $6$ or $3$.
Neither is a power of $2$ modulo $7$. Consequently
\[
\boxed{N<49\quad\Longrightarrow\quad\gcd(M,N)=1.}
\]
Indeed, for $35<N<49$, a prime factor of a nontrivial $g$ would have
to be $7$, since $N\ge5g$; no multiple of $7$ in this interval is
odd and coprime to $3$.

\subsection{Using congruences modulo eight and five}
\label{sec:mod-eight-families}
If $7\mid g$, reducing the two partitions modulo $7$ gives
\[
\begin{array}{c|c|c|c|c}
q\bmod7&m_2\bmod6&m_1\bmod6&n_1\bmod6&N\bmod6\\\hline
3&4&2&4&5\\
5&4&2&2&1
\end{array}
\]
Indeed $p\equiv1\pmod7$ requires
$2^{m_1}\equiv q^{n_1}\equiv4\pmod7$. Substituting the possible
residues into $2^{m_2}+q^{n_1+N}\equiv1$, with $\gcd(N,6)=1$,
leaves the two rows. In particular $M\equiv4\pmod6$ and
$n_2\equiv3\pmod6$.

\begin{lemma}\label{lem:q-one-mod-eight}
In the remaining case, $q\equiv1\pmod{48}$ and $m_2$ is even and
at least $4$.
\end{lemma}
\begin{proof}
Suppose $q\equiv5\pmod8$. Then $m_2=2$ and
$v_2(\operatorname{ord}_q(2))=2$, so $4\mid M$.
Thus $5\mid A$ and $q^N\equiv1\pmod5$. Since $N$ is odd,
$q\equiv1\pmod5$. But $m_1=M+2\equiv2\pmod4$, giving
$p=2^{m_1}+q^{n_1}\equiv4+1\equiv0\pmod5$, impossible.
The remaining class $q\equiv1\pmod8$, together with even
$m_2=v_2(q-1)$, gives $q\equiv1\pmod{16}$.
Combine this with $q\equiv1\pmod3$.
\end{proof}

\begin{lemma}\label{lem:mod-five-split}
The congruences modulo $5$ split as follows:
\[
\begin{array}{c|c|c}
 &q\bmod5&\text{further conditions}\\\hline
4\mid M&1&m_2\equiv0\pmod4,\quad5\nmid N,\quad p\equiv2\pmod5\\
M\equiv2\pmod4&2,3&q^N\equiv3\pmod5,\quad m_2+n_1\equiv2\pmod4
\end{array}
\]
Consequently
\[
q\equiv1,97,193\pmod{240},\qquad
q^2\equiv1,49\pmod{240}.
\]
\end{lemma}
\begin{proof}
When $4\mid M$, we have $q\equiv1\pmod5$, as above.
Then $p\equiv2^{m_2}+1\pmod5$ excludes $m_2\equiv2\pmod4$.
The assertion $5\nmid N$ follows from \eqref{eq:ordinary-q-minus-one}.
When $M\equiv2\pmod4$, put $P=2^{m_2}$ and $Q=q^{n_1}$.
Both are $\pm1$ modulo $5$, and $p\equiv Q-P$ cannot vanish.
Hence $Q\equiv-P\pmod5$. Since $AQ\equiv2^M-1\equiv3$,
\[
q^N=AP+1\equiv-3+1\equiv3\pmod5.
\]
Thus $q\equiv2$ or $3\pmod5$. The relation $Q\equiv-P$ gives
$m_2+n_1\equiv2\pmod4$. Combining the three possible residues
of $q$ modulo $5$ with $q\equiv1\pmod{48}$ gives the last assertions.
\end{proof}
In the second row, $q\equiv2\pmod5$ requires $N\equiv3\pmod4$,
and $q\equiv3\pmod5$ requires $N\equiv1\pmod4$.
Also
\[
p\equiv65\pmod{96},\qquad p\equiv257\text{ or }353\pmod{480}.
\]
Indeed $q^{n_1}\equiv1\pmod{96}$ and
$2^{m_1}\equiv64\pmod{96}$, while $p\equiv2,3\pmod5$.

% New deductions from the common-divisor residue calculations.
\begin{lemma}\label{lem:exclude-seventeen}
$17\nmid\gcd(M,N)$. Consequently gaps are coprime whenever
$N=5^a17^b$, with $a,b\ge0$.
\end{lemma}
\begin{proof}
Suppose $17\mid g$. The nonzero squares modulo $17$ are
$1,2,4,8,9,13,15,16$, and
$\langle4\rangle=\{1,4,16,13\}$.
In $2^{m_1}+q^{n_1}\equiv1$, only
$2^{m_1}\equiv16$, $q^{n_1}\equiv2$ are possible.
Thus $m_1\equiv4\pmod8$. Since $\operatorname{ord}_{17}(2)=8$
and $n_1$ is even, $q$ has order $16$ and $n_1\equiv2\pmod4$.
The odd power $q^{n_2}$ also has order $16$.
In $q^{n_2}\equiv1-2^{m_2}$, the choices
$2^{m_2}=1,16$ give $0,2$, so only $4,13$ remain.
Hence $m_2\equiv2\pmod4$ and $M\equiv2\pmod4$.
But Lemma~\ref{lem:mod-five-split} then requires
$m_2+n_1\equiv2\pmod4$, a contradiction.
\end{proof}

\subsection{Sharpening the height and checking residue families}
\begin{lemma}\label{lem:sharp-h}
In the remaining case,
\[
h<q^2,\qquad (2h-1)n_1<N,\qquad
h<\frac{N+2}{4}.
\]
In particular $(2g-1)n_1<N$ when $g>1$.
\end{lemma}
\begin{proof}
The classes modulo $240$ give $q\ge97$. Put $\lambda=\log q$ and
$\beta=\log(27/26)$. Suppose $h\ge q^2$.
Since $M=hE$ and $E\log2\ge2\log(q^2-1)>3\lambda$,
$M>3h\lambda/\log2$.
In \eqref{eq:baker-additional} replace $1-1/(2h)$ by $0.99$.
The function
\[
\frac{0.99x\log2-\beta}
 {\max\{\log(3x/(2\lambda))+0.38,10\}^2}
\]
is increasing for $x\ge3h\lambda/\log2$. Substitution gives
\[
2.97h\lambda-\beta
 <25.2\lambda^2(\log h+3)^2.
\]
The ratio $(2.97h\lambda-\beta)/(\log h+3)^2$ also increases
for $h\ge q^2$. Hence
\[
q^2<9\lambda(2\lambda+3)^2.
\]
At $q=97$ the right side is less than $6078<9409$;
$q^2/[\log q(2\log q+3)^2]$ increases for $q\ge97$.
This proves $h<q^2$.

Write $e=n_1$, $c=m_2$, $Q=q^e$. Since $c\ge4$, $e\ge2$,
\[
2h\log\frac{Q}{Q-1}
 <\frac{2q^2}{q^2-1}<4\log2\le c\log2.
\]
Using $2^E\ge(Q-1)^2$ and $2^{m_1}<q^{n_2}$, we obtain
\[
q^{N+e}>2^{c+hE}\ge2^c(Q-1)^{2h}>Q^{2h}.
\]
Thus $(2h-1)e<N$. Now use $e\ge2$ and $g\mid h$.
\end{proof}

\subsection{A uniform bound on the odd gap}
\begin{lemma}\label{lem:uniform-N}
In the remaining case $N\le2519$, $n_2\le5037$, $h\le630$, and
$M<5037\log_2 q$.
\end{lemma}
\begin{proof}
Use Laurent's Corollary~2 with $D=1$, $m=10$, $C_2=25.2$.
Its hypothesis is
$H_i=\log A_i\ge\max\{\mathrm h(\alpha_i),|\log\alpha_i|,1\}$
for each $i$ separately \cite[p.~328]{Laurent}. For
$(\alpha_1,\alpha_2)=(2,q)$ take $H_1=1$ and $H_2=\log q$.
The bases are multiplicatively independent. Since
$2^{m_1}<q^{n_2}$ and $n_2=N+\rho<2N$ by Lemma~\ref{lem:sharp-h},
\[
 b'=\frac{m_1}{\log q}+n_2
 <(1+1/\log2)n_2<5N.
\]
Applying the corollary to $n_2\log q-m_1\log2$, then
\eqref{eq:lambda-upper}, and dividing by $\log q$ gives
\[
 N <25.2\max\{10,\log(5N)+0.38\}^2+1.
\]
The function $(x-1)/\max\{10,\log(5x)+0.38\}^2$ is increasing
for $x\ge2521$. On the logarithmic branch this follows from
$\log(5x)+0.38>2(1-1/x)$. At $x=2521$ the maximum is $10$
and the ratio is $25.2$. Hence $N<2521$. Now $N$ is odd,
$\rho<N$, and $h<(N+2)/4$. The bound on $M$ follows from
$M<m_1<n_2\log_2q$.
\end{proof}

\paragraph{Finiteness as $q$ varies.}
Put $n=n_2$ and write $m_1=jn+s$, $0\le s<n$. Then
$\alpha=2^{s/n}$, $y=2^j$ and $x=y\alpha$ satisfy
\[
 0<q^n-x^n=q^\rho-2^c<q^\rho,
 \qquad 0<q-x<q^{1-N}<1.
\]
Indeed the quotient of $q^n-x^n$ by $q-x$ exceeds $q^{n-1}$.
Thus $s\ne0$, and $\alpha$ is irrational by its $2$-valuation
(or by taking the $n$th power of a rational in lowest terms).
As $q>y$ this gives
\[
 0<q/y-\alpha<y^{-N}.
\]
The fractions are reduced. There are only finitely many $(n,s)$
by Lemma~\ref{lem:uniform-N}, and $N\ge5$. Roth's theorem
\cite{Roth1955} gives finitely many such fractions for each
$\alpha$ (the case $y=1$ is immediate from $q-x<1$).
Hence the remaining system has finitely many solutions over all $q$.
This proof supplies no computable last denominator. An effective
inequality $|\alpha-q/y|>C y^{-\mu}$ with explicit $C>0$ and
$\mu<N$ would give $y<C^{-1/(N-\mu)}$ for that exponent family.

\begin{lemma}\label{lem:cubic-large-exponent}
If $3\mid n_2$, then $n_2<6\rho$, $h\le2$, and $\gcd(M,N)=1$.
Consequently $7\nmid\gcd(M,N)$ for all remaining solutions.
\end{lemma}
\begin{proof}
Write $n_2=3k$, $Y=q^k$, $m_1=3j+s$ and $X=2^{m_1/3}$.
Then $0<Y-X<q^\rho/Y^2$. Since $\rho<n_2/2$,
if $s=0$ the distinct integers $Y,X$ would give
$Y^3-X^3>Y^2=q^{2k}>q^\rho$.
Write $T=2^j$ and $\theta=\sqrt[3]2$. For $s=1$, apply Bennett's
$|\theta-P/D|>1/(4D^{5/2})$ \cite[eq.~(4)]{BennettCubic}
to $P/D=Y/T$. This gives
$q^\rho>Y^2/(4T^{3/2})>Y^{1/2}/4$.
For $s=2$, use $P/D=2T/Y$ and $\theta X=2T$. Then
$0<\theta-2T/Y<\theta q^\rho/Y^3$, so
$q^\rho>Y^{1/2}/(4\theta)$.
In either case $q^{k/2-\rho}<4\theta<\sqrt{97}\le\sqrt q$.
Thus $k\le2\rho$ and, since $k$ is odd, $k<2\rho$.
Now $2h\rho<n_2$ gives $h<3$. Since $g\mid h$ and $g$ is odd,
$g=1$. Finally the modulo-$7$ table above forces $3\mid n_2$
whenever $7\mid g$.
\end{proof}
In particular $N<65$ implies $g=1$. Otherwise $N\ge5g$ and all
prime factors of $g$ are at least $11$. Below $65$ this leaves
$g=11$ and $N=55$, for which $N=5g$ gives $\rho=2$ and $n_2=57$,
contrary to Lemma~\ref{lem:cubic-large-exponent}. More generally
$N=5g$, $g>1$, requires $g\equiv1\pmod3$, since $n_2=5g+2$.
Also $g\le\lfloor2519/5\rfloor=503$. Since $11^3>503$, a
nontrivial $g$ is a prime or a product of two primes counted
with multiplicity.

\begin{lemma}\label{lem:aurifeuillean-height}
Suppose $r\equiv4\pmod8$. If
\[
\frac{8h}{q^{n_1/2}-1}\le m_2\log2,
\]
then $(4h-1)n_1<N$. The hypothesis holds whenever $h\le32$.

\end{lemma}
\begin{proof}
Write $r=8a+4$. Since $q^{n_1}\mid2^r-1$ and
$\operatorname{ord}_q(2)=r$, the full power $q^{n_1}$ divides
\[
2^{r/2}+1
=(2^{2a+1}-2^{a+1}+1)(2^{2a+1}+2^{a+1}+1).
\]
The two factors are odd and their difference is a power of $2$,
so they are coprime. Hence
\[
q^{n_1}\le2^{r/4}+2^{(r+4)/8}+1
 <(2^{r/8}+1)^2,
\qquad 2^r>(q^{n_1/2}-1)^8.
\]
Here $E=r$. Put $e=n_1$, $c=m_2$. The hypothesis and
$\log(1+x)<x$ give
\[
8h\log\frac{q^{e/2}}{q^{e/2}-1}<c\log2.
\]
Therefore
\[
q^{N+e}>2^{c+hr}
 >2^c(q^{e/2}-1)^{8h}>q^{4he}.
\]
For $h\le32$, use $q\ge97$, $e\ge2$, $c\ge4$:
$8h/(q^{e/2}-1)\le8/3<4\log2$.

\end{proof}
Pomerance \cite[\S2, (1)]{PomCycl} records the corresponding
Aurifeuillean factorization for $\Phi_r(2)$ when $r\equiv4\pmod8$.
Brent \cite[\S\S2--3]{Brent} develops its relation to grouping
roots of unity.

For a fixed prime $d\mid g$, retain the even values $e=n_1$ and
multiples $h$ of $d$ satisfying $(2h-1)e<N$.
For each, reduce
\[
2^{m_1}=1-q^e,\qquad 2^{m_2}=1-q^{N+e}\pmod d
\]
with $m_1,m_2,E$ even, $m_1-m_2=hE$, and the modulo-$5$ split.
This is a finite residue calculation; no upper bound on $q$ is used.
The script \texttt{check\_common\_gap\_residues.lua} gives:
\[
\begin{array}{c|c|c|c|l}
N&d&n_1&h&q\bmod d\\\hline
49&7&2&7&5\\
55&11&2&11&3,5,6\\
65&13&2&13&2\\
77&7&4&7&3\\
77&11&2&11&5,8\\
85&17&\multicolumn{3}{c}{\text{excluded}}\\
91&7&2&7,14&5\\
91&13&2&13&2,6\\
121&11&2&11,22&6
\end{array}
\]
These rows record the residue test alone. Lemma~\ref{lem:cubic-large-exponent}
now excludes all the rows with $d=7$ and the row $N=55$, $d=11$.
For $N=5d$, the same test proves coprimality for
\[
d\in\{7,17,29,41,43,73,89,109,113\}.
\]
Here $g>1$ would force $g=d$, $e=2$, $h=d$.
Since $5\mid N$, the modulo-$5$ split gives
$m_2\equiv0$, $m_1\equiv2\pmod4$.

An algebraic version of this last test uses $t=\operatorname{ord}_d(16)$.
Compute over $\mathbb F_d$ the gcd of
\[
X^{d-1}-1,\qquad (1-X^2)^t-4^t,\qquad (1-X^7)^t-1.
\]
If it is $1$, the family $N=5d$ is excluded. The exponent $7$ comes
from $N+2=5d+2\equiv7\pmod{d-1}$.
More generally use the two cosets $\langle16\rangle$ and
$4\langle16\rangle$, retaining the permitted classes of $m_1,m_2$
modulo $4$. A general extension would need to show that all the
residue sets surviving the height bound are empty, possibly after
combining further moduli. No such uniform assertion is proved here.

\subsection{Forcing a common factor}
\label{sec:forcing-common-factor}
Let $\ell=jM+1>q$ be a primitive Mersenne factor at exponent $M$.
If $\gcd(j,N)=1$, then
\[
1<\operatorname{ord}_\ell(q)\mid\gcd(M,N).
\]
Indeed, $\ell\mid A$ gives $\operatorname{ord}_\ell(q)\mid N$,
while the order also divides $\ell-1=jM$. Coprimality with $j$
forces it to divide $M$, and $1<q<\ell$ excludes order $1$.
Thus such a factor contradicts $N<65$. A multiplier $j$ supported
on $2,3$ automatically passes the coprimality test; if $4\mid M$,
then $5\nmid N$, so factors of $5$ in $j$ are allowed as well.
The missing general step is to force a primitive factor with this
multiplier condition and $\ell>q$. Zsigmondy's theorem alone does
not supply either condition.

The condition $\ell>q$ may be weakened to $\ell\nmid q-1$.
There is also a product version. Put $u=(q-1)/2^{m_2}$ and
\[
D=\prod_{\substack{\ell\mid A,\ \operatorname{ord}_\ell(2)=M\\
                    \gcd((\ell-1)/M,N)=1}}
       \ell^{v_\ell(A)}.
\]
If $g=1$, every prime in this product divides $q-1$.
Moreover $\ell>M>N$, so lifting gives
$v_\ell(A)=v_\ell(q-1)$, and therefore
\[
\boxed{g=1\ \Longrightarrow\ D\mid u.}
\]
Thus $D\nmid u$, in particular $D>u$, forces $g>1$.
This allows several smaller primitive factors to replace one factor
larger than $q$. The prime $q$ is excluded from the product even when
$\operatorname{ord}_q(2)=M$.

The same test also applies inside the common-divisor branch.
Since $\gcd(E,N)=1$, form $D_E$ from the full prime-power factors
of $(2^E-1)/q^{n_1}$ satisfying
\[
\gcd\left(\frac{\ell-1}{\operatorname{ord}_\ell(2)},N\right)=1.
\]
Each such $\ell$ divides $q-1$, and lifting gives
\[
\boxed{D_E\mid uN.}
\]
Thus a lower bound $D_E>uN$ would exclude that family without
first proving the gaps coprime.

For $j=2$, necessarily $4\mid M$: otherwise $\ell=2M+1\equiv5\pmod8$
and $2^M\equiv-1\pmod\ell$.
Even with $4\mid M$, primality of $2M+1$ and order exactly $M$
still require proof.

The multiplier $j=(\ell-1)/\operatorname{ord}_\ell(2)$ is the
\emph{residual index} of $2$ modulo $\ell$.
For a prime $s\mid N$, the condition $s\mid j$ is equivalent to
$\ell$ splitting completely in $\mathbb Q(\zeta_s,2^{1/s})$.
Here $\ell>M>N$, so no ramified primes intervene.
This identifies the desired index restriction with avoidance of a
finite collection of splitting conditions. A useful general result
would bound the part of the primitive factorization supported on
those splitting primes, strongly enough to force $D>u$.

\subsection{A finite target when q is bounded in terms of d}
% Added from the question about whether large common divisors are prohibitive.
\begin{lemma}
If a prime $d\mid\gcd(M,N)$ satisfies $q\le d^2$, then $d<10^5$.
\end{lemma}
\begin{proof}
Suppose $d\ge10^5$, and put $t=\log d$, $\lambda=\log q\le2t$.
Write $H_\lambda(x)=\max\{\log(3x/\lambda)+0.38,10\}$.
The function
\[
\frac{x\log2/2-\log(27/26)}{H_\lambda(x)^2}
\]
is increasing for the positive $x$ used below. Also
$\lambda H_\lambda(x)$ is increasing in $\lambda$.
These allow the lower bounds on $M$ to be substituted into
\eqref{eq:baker-M}.

We have $dE\mid M$, so
$M\ge2d\log_2(q^2-1)>3d\lambda/\log2$.
At $x=3d\lambda/\log2$, $H_\lambda(x)<t+3$.
Consequently \eqref{eq:baker-M} implies
\[
d<34t(t+3)^2.
\]
At $d=10^5$ the right side is less than $83000$; moreover
$d/[\log d(\log d+3)^2]$ increases thereafter, a contradiction.
\end{proof}
Thus $d\ge10^5$ requires $q>d^2$. More generally, imposing
$q\le d^K$ for any fixed $K>0$ gives an effective finite bound
on $d$ by the same argument: the necessary bound has the form
$d=O_K((\log d)^3)$.
This is a finite target for a specified family, not an absolute bound
on common divisors when $q$ is unrestricted.

\section{Factorization results and their use here}
\label{sec:factor-results}
% Original notes Sections 8.2, 8.4, 9.5--9.6 restored; new consequences marked.
For $u=(q-1)/2^{m_2}$, the original factorization $A=u\Psi$ gives
\[
\operatorname{lcm}\left(\operatorname{ord}_q(2),
 \{\operatorname{ord}_\ell(2):\ell\mid u\text{ prime}\}\right)\mid M.
\]
Also $M/r<2N$ for $r=\operatorname{ord}_q(2)$, since
$2^{M/2}<q^N<2^{rN}$.
Zsigmondy's theorem applied to $q^N-1$, with odd $N\geq3$, has no
exception. It supplies a prime $\ell_0\mid\Psi$ such that
\[
\operatorname{ord}_{\ell_0}(q)=N,\qquad
\ell_0\equiv1\pmod{2N},\qquad\operatorname{ord}_{\ell_0}(2)\mid M.
\]
Conversely, a primitive Mersenne factor at $M\neq6$ has order $M$
and lies among the primes dividing $q(q^N-1)$. Thus fixing $(q,N)$
restricts $M$ to the orders at base $2$ of those primes, with $6$
as the primitive-divisor exception.

\subsection{Gcd identities and primitive-factor counts}
Cambraia et al.\ \cite[\S2, Proposition 9]{Cambraia} collect the
same-base gcd and valuation facts used here. In particular,
\[
\gcd(2^a-1,2^b-1)=2^{\gcd(a,b)}-1,
\]
\[
\gcd\left(2^a-1,\frac{2^{ab}-1}{2^a-1}\right)=\gcd(2^a-1,b).
\]
The second follows by reducing the geometric sum modulo $2^a-1$.
Their Proposition 2 uses coprime exponents and a primitive factor at
the product exponent to force an additional prime.

Write $\omega(n)$ for the number of distinct prime factors and
$\tau(n)$ for the number of positive divisors. The primitive-factor
argument in Kova\v{c}--Luca \cite[\S3]{KL} gives, more precisely,
\[
\omega(2^M-1)\geq\tau(M)-1-\mathbf1_{6\mid M},\qquad
\omega(A)\geq\tau(M)-2-\mathbf1_{6\mid M}.
\]
Choose a primitive prime at every divisor of $M$ except $1,6$;
different orders give different primes. Removing $q$ loses one prime.

\begin{lemma}
For the remaining case $A>1$, $q>3$,
\[
\omega(2^M-1)\geq4,\qquad\omega(A)\geq3.
\]
\end{lemma}
\begin{proof}
Suppose $\omega(2^M-1)\leq3$. The preceding count, with $M$ even,
leaves $M=2,4,6,8$ or $M=2s$ for an odd prime $s$.
The size bound excludes $2,4$; $2^6-1=3^2\cdot7$ and
$2^8-1=3\cdot5\cdot17$ exclude $6,8$ since $q>3$ and $n_1\geq2$.
For $s\geq5$, the three primes are $3$, a primitive factor at $s$,
and a primitive factor at $2s$. Hence
\[
2^s-1=P^f,\qquad2^s+1=3Q^e
\]
for distinct primes $P,Q$, using $v_3(2^s+1)=1$.
A Mersenne number $2^s-1$ is not a nontrivial perfect power:
an even exponent is excluded modulo $4$, and an odd exponent greater
than $1$ makes $P^f+1$ have an odd factor greater than $1$.
Thus $f=1$, so $q=Q$ and $e=n_1$ is even.
But $2^s+1=3Q^e$ is then impossible modulo $8$.
\end{proof}
% This consequence is derived here; the low-factor classification also
% appears in Cambraia et al., Theorems 14--15.

\subsection{Transferring factors into C}
For a prime $s\mid\gcd(M,N)$, put $a=M/s$, $b=N/s$,
$Y=2^a q^b$ and $C=A/(Y-1)$. Then
\[
Cp=\Phi_s(Y),\qquad C<\frac{Y^{s/2}}{2q(Y-1)},\qquad
\gcd(C,Y-1)\mid s.
\]
For any divisor $D\mid A$,
\[
\frac{D}{\gcd(D,Y-1)}\mid C.
\]
Taking $D=(2^a-1)/q^{v_q(2^a-1)}$ gives
\[
\boxed{\frac{D}{\gcd(D,q^b-1)}\mid C,}
\]
since $Y-1\equiv q^b-1\pmod D$.
Thus a useful lower bound on $C$ requires control of a mixed-base gcd.
For $\ell\mid C$, $\ell\neq s$, we have $\operatorname{ord}_\ell(Y)=s$;
therefore, with $r_\ell=\operatorname{ord}_\ell(2)$ and
$t_\ell=\operatorname{ord}_\ell(q)$,
\[
v_s(r_\ell)=v_s(M)\quad\text{or}\quad v_s(t_\ell)=v_s(N).
\]
If both valuations were smaller, both factors of $Y$ would be $1$
modulo $\ell$. The converse need not hold because of cancellation.

\subsection{Primitive factors on the two sides}
% Added from the discussion of Zsigmondy's theorem and the full odd-part equality.
Let $d\ge5$ be a prime dividing both gaps, $U=2^{M/d}$,
$V=q^{N/d}$, $B=UV-1$, and $C=A/B$. The full equality is
\[
2^M-1=q^{n_1}\operatorname{odd}(q^N-1),\qquad
\operatorname{odd}(q^N-1)=A.
\]
Thus every prime power in $2^M-1$ outside the base $q$ occurs
with exactly the same valuation in $q^N-1$.

Zsigmondy's theorem \cite[Theorem 13]{GyorySmyth} supplies a prime
of order $M$ at base $2$ and a prime of order $N$ at base $q$.
Here $M\ge10$ and $N\ge5$ is odd, so there are no exceptions.
The former prime belongs to $A$ unless it is $q$ itself; the latter
always belongs to $A$. Such primes satisfy respectively
$\ell\equiv1\pmod M$ and $\ell\equiv1\pmod{2N}$.
If the same prime has both full orders, then
$\ell\equiv1\pmod{\operatorname{lcm}(M,N)}$.

For $\ell\mid A$, $\ell\ne d$, put
$\alpha=\operatorname{ord}_\ell(U)$ and
$\beta=\operatorname{ord}_\ell(V)$. Both lie in $\{1,d\}$.
The factor allocation is
\[
\begin{array}{c|c|c}
\alpha&\beta&\text{location of the full }\ell\text{-part of }A\\\hline
1&1&B\\
1&d&C\\
d&1&C\\
d&d&B\text{ if }UV\equiv1\pmod\ell;\ C\text{ otherwise}.
\end{array}
\]
Indeed $Cp=\Phi_d(UV)$ and $\gcd(B,C)\mid d$.
In the last row, $U,V$ lie in the same subgroup of order $d$;
their product can be $1$. Primitivity at two different bases does
not exclude this cancellation.

\begin{lemma}
Put $G=\gcd(U-1,V-1)$. Then $B=GH$ with
\[
\gcd(G,H)=1,\qquad H>U,
\]
and every prime $\ell\mid H$ satisfies
$\operatorname{ord}_\ell(U)=\operatorname{ord}_\ell(V)=d$.
\end{lemma}
\begin{proof}
For $\ell\mid B$, the residues $U,V$ are inverse, so their orders
are equal. If this order is $1$ and $\ell\ne d$, lifting and the
full gap equations give
\[
v_\ell(B)=v_\ell(A)=v_\ell(U-1)=v_\ell(V-1).
\]
If $d\mid B$, then $U\equiv V\equiv1\pmod d$ and
$v_d(\Phi_d(UV))=1$. Since $p>d$, $v_d(C)=1$.
Lifting again gives
$v_d(B)=v_d(U-1)=v_d(V-1)$.
Thus the entire part of $B$ whose primes have order $1$ is exactly
$G$, with no such primes left in $H$. Finally $G\le V-1$ gives
$H\ge(UV-1)/(V-1)>U$.
\end{proof}

For a primitive Mersenne factor $\ell\ne q$ at exponent $M$,
$\alpha=d$. Writing $t=\operatorname{ord}_\ell(q)$, if
$v_d(t)<v_d(N)$ then $\beta=1$, so the full $\ell$-part of $A$
is forced into $C$. Symmetrically, a primitive factor at exponent
$N$ for base $q$ is forced into $C$ if
$v_d(\operatorname{ord}_\ell(2))<v_d(M)$.
Products of distinct forced prime powers must fit below
\[
C<\frac{(UV)^{d/2}}{2^{m_2/2}q^{n_1/2}(UV-1)}.
\]
The remaining target is to force enough such factors into $C$, or
to bound the simultaneous primitive part $H$ below its required
size. Zsigmondy's existence assertion alone does neither.

\subsection{Reading and limits of the cited results}
Pomerance \cite[\S2]{PomCycl} supplies the primitive/intrinsic
cyclotomic decomposition, bounds for $\Phi_d(2)$, and the Aurifeuilian
splitting for $d\equiv4\pmod8$. These concern individual factors.
Cambraia et al.\ \cite[\S\S2--4]{Cambraia} give the gcd identities,
valuation formula, and classification when $\omega(2^M-1)\leq3$;
the preceding lemma specializes that classification to our equations.
Their final section studies a different sum of Mersenne numbers.

Kova\v{c}--Luca \cite{KL} study sums of divisor counts over exponents;
their pointwise primitive-factor bound is what we use above.
Erd\H{o}s--Kiss--Pomerance \cite{EKP} study reciprocal sums of prime
divisors, including averages and consecutive exponents. Neither result
currently controls the particular mixed-base gcd needed for $C$.
Pomerance's earlier paper \cite{PomPrim} treats reciprocal sums over
divisors of prescribed order.

For the proposed descent, compare Bilu \cite[\S2, Proposition 2.1 and
Corollary 2.2]{Bilu}: coprimality and a perfect-power product force powers
of ideals. The identity $\Phi_s(Y)=Cp$ does not provide that
perfect-power hypothesis. A transfer of the argument would need to
supply it.

\subsection{Cyclotomic structure and further reading}
% Further reading and open targets from the current discussion.
For a primitive $M$-th root of unity $\zeta_M$,
\[
\Phi_M(2)=\operatorname{Norm}_{\mathbb Q(\zeta_M)/\mathbb Q}(2-\zeta_M).
\]
A primitive Mersenne factor $\ell$ corresponds to a prime ideal
where $\zeta_M$ reduces to $2$. Subgroups of the Galois group group
the conjugates into relative norms. This can give algebraic
factorizations and, under suitable specialization, integer factors.
Brent \cite[\S\S2.2--2.3,3,4.4]{Brent} develops the Gauss and
Aurifeuillean identities from these root groupings and gives algorithms
for computing their factors. Lemma~\ref{lem:aurifeuillean-height}
is one use here: the whole $q^{n_1}$ must fit into one coprime factor.
For the field structure, use Washington \cite[Chapters 2--3]{Washington};
Chapter 6 gives the Gauss-sum background, and Chapter 8 treats
cyclotomic units. A further target is to compute relative norms in
small subfields and determine whether they separate the part containing
$q^{n_1}$ from factors forced into $A$.

The magnitude condition $\ell>q$ remains a separate route.
Stewart \cite[Theorem 1.1, \S5]{Stewart} proves, for all sufficiently
large $M$, an effective bound
\[
P(\Phi_M(2))>
M\exp\!\left(\frac{\log M}{104\log\log M}\right),
\]
where $P$ denotes the largest prime factor. Such a factor is primitive.
The proof combines the size of the cyclotomic value with bounds on
prime-adic multiplicities. To use it here, one must compare its bound
with $q$ in the family under consideration and control the multiplier
of the resulting prime. The current general bound on $q$ in terms
of $M$ is exponential, so this does not yet establish $\ell>q$ throughout.

The alternative $\ell\nmid q-1$ asks for control of shared factors
with $q-1$. The products $D,D_E$ above make the required divisibility
or size statement precise; logical weakness alone does not make
that statement easier to prove.
For the residual index $j$, Moree \cite[\S\S1,5,6]{MoreeNear}
provides a complementary viewpoint.
The condition $\operatorname{ord}_\ell(2)=(\ell-1)/j$ is precisely
his near-primitive-root problem. Lenstra--Moree--Stevenhagen
\cite[\S\S2,5,6]{LMS} express the compatibility of the splitting
conditions using Galois groups and character sums. Their density
results assume GRH and vary $\ell$; they do not force a suitable
prime divisor of a specified $\Phi_M(2)$. Here the additional target
is the bound $D>(q-1)/2^{m_2}$ from
Section~\ref{sec:forcing-common-factor}.

For a different route to a bound on $N$, Corvaja's short survey
\cite[\S1]{CorvajaSurvey} explains the Subspace Theorem bounds for
gcds of powers and the extension to a fixed finitely generated
multiplicative group. In our case
\[
A=\gcd(2^M-1,q^N-1)>3q^{N-1}.
\]
The strict inequality follows from $2^{m_2}\le(q-1)/3$.
An effective bound, valid for our family and uniform in $q$, of the form
\[
\log\gcd(2^M-1,q^N-1)
 \le\theta N\log q+K\log q,\qquad \theta<1,
\]
with fixed $\theta,K$, would give $N<(K+1)/(1-\theta)$.
This is a target, not a quoted theorem. Fixed-group results do not
supply uniformity as $q$ varies, and their ineffectivity also matters
for a finite computation. The survey connects this approach with
complex analysis and the support problem.

For computational algebraic number theory, Alvarado et al.
\cite[\S\S3--6]{SUnitAlgorithms} give explicit archimedean and
$\ell$-adic bounds, lattice reduction, and sieves for $S$-unit equations
over number fields. A use here would first require a reduction to
controlled fields and prime supports. The original equality gives a
four-term unit relation over $S=\{2,q\}$, not directly their
two-variable equation $x+y=1$. In a cyclotomic-field approach the
uncontrolled primes in $A$ remain the obstacle to fixing that support.

\section{An order obstruction from power reciprocity}
\label{sec:reciprocity}

\subsection{The order test}

Work in the remaining case $q>3$, $A>1$ and $q\nmid M$. Keep
$r=\operatorname{ord}_q(2)$, $\rho=v_q(2^r-1)=n_1$ even, and
$E=\operatorname{lcm}(2,r)$. Then
\[
a=\frac{2^E-1}{q^\rho}\mid A,\qquad \gcd(E,N)=1.
\]
Indeed, $r$ divides both $M$ and $q-1$, so
$\gcd(M,N,q-1)=1$ from Section~\ref{sec:p-common-gap} gives
$\gcd(r,N)=1$. Also $N$ is odd.
Since $A2^{m_2}=q^N-1$, the order of $q$ modulo any prime power
$\ell^f\mid a$ divides $N$. To exclude $q$ for all $M,N$, it is enough
to find one such order with a prime divisor in common with $E$.

Fix a prime power $m=s^b$ dividing $r$. If $\ell\mid a$ is prime and
$m\mid\ell-1$, then
\[
q^{(\ell-1)/m}\not\equiv1\pmod\ell
\quad\Longrightarrow\quad s\mid\operatorname{ord}_{\ell}(q).
\]
An order prime to $s$ would divide $(\ell-1)/m$, since we only
removed factors of $s$ from $\ell-1$.

For a solution, $q$ must therefore be an $m$-th power at each such
prime $\ell$. The power-residue symbol records this condition as an
$m$-th root of unity, with value $1$ when $q$ is a power. We choose a
cyclotomic divisor of $2^E-1$ and multiply the symbols at its prime
ideal factors away from $s$ and $q$, with multiplicity. If the
product differs from $1$, at least one of those primes gives the
required obstruction.

The reciprocity law here is the product formula for local
norm-residue symbols. It relates the product at the unknown factors
to local calculations at the known prime $q$, once the contribution
at $s$ is accounted for. This gives an existence test for the prime
$\ell$ without first finding it. The aim is to force an obstruction
uniformly for the remaining $q$, or to prove an effective bound on
the exceptions to this test.

\subsection{Choosing the cyclotomic factor}

To select primes for the test, set $B=2^{E/m}$ and take the
cyclotomic value
\[
H_s=\Phi_m(B)
  =\frac{2^E-1}{2^{E/s}-1},
\qquad a_s=\frac{H_s}{q^\rho}.
\]
Since $m\mid r$, the order of $B$ modulo $q$ is
$r/\gcd(r,E/m)=m$. The equality uses $\gcd(m,E/r)=1$,
as $E/r=2$ only when $r$ is odd.
In particular, $2^{E/s}=B^{m/s}\not\equiv1\pmod q$. Thus $H_s$
contains the full $q$-part of $2^E-1$ and $a_s\mid a$.

Work in $K=\mathbb Q(\zeta_m)$ with $\zeta_m$ a primitive $m$-th
root of unity. Its ring of integers is $\mathcal O_K=\mathbb Z[\zeta_m]$
\cite[Chapter 2]{Washington}. Set $\beta=B-\zeta_m$, whose norm to
$\mathbb Q$ is $H_s$. Let $P$ denote a nonzero prime ideal of
$\mathcal O_K$. For nonzero $\alpha\in\mathcal O_K$, the valuation
$v_P(\alpha)$ is the exponent of $P$ in the factorization of
$(\alpha)$. In particular,
\[
(\beta)=\beta\mathcal O_K=\prod_{P\mid(\beta)}P^{v_P(\beta)},
\qquad
J=\prod_{\substack{P\mid(\beta)\\ P\nmid(sq)}}P^{v_P(\beta)}.
\]
Thus $J$ keeps the prime ideal factors of $(\beta)$ away from $s$
and $q$, with their multiplicities. In terms of containment,
$(\beta)\subseteq J$, and $J\subseteq P$ for each prime ideal $P\mid J$.

Fix a prime ideal $P\mid J$ and let $\ell$ be the rational prime
below it, so $P\cap\mathbb Z=\ell\mathbb Z$. Reduction at $P$
is the quotient map $\mathcal O_K\to\mathcal O_K/P$.
The quotient $\mathcal O_K/P$ is the residue field at $P$.
It is a field because $P$ is maximal.
Here $\beta\in P$, so $\zeta_m$ maps to $B\bmod\ell$. Since
$\mathcal O_K=\mathbb Z[\zeta_m]$, the image is generated by
$\mathbb F_\ell$ and this residue of $\zeta_m$, which already lies
in $\mathbb F_\ell$. Hence
\[
\mathcal O_K/P\simeq\mathbb F_\ell.
\]
Since $\ell\nmid m$, the reduction of $\zeta_m$ still has order $m$.
Indeed, $X^m-1$ has distinct roots modulo $\ell$, so its factor
$\Phi_m(X)$ has no common root with any $\Phi_d(X)$ for a proper
divisor $d$ of $m$. Therefore $B\bmod\ell$ has order $m$ and
$m\mid\ell-1$.

The $m$-th power-residue symbol $(q/P)_m$ is the root of unity
that reduces to $q^{(\ell-1)/m}$ at $P$. Equivalently, it is the
factor by which Frobenius multiplies an $m$-th root of $q$ in the
Kummer extension of $K$ \cite[\S\S2,4]{StevenhagenKummer}.
Multiplying with the multiplicities in $J$ gives
\[
(q/J)_m:=\prod_{P\mid J}(q/P)_m^{v_P(\beta)}.
\]
If $(q/J)_m\ne1$, some symbol $(q/P)_m\ne1$. The order test
then gives $s\mid\operatorname{ord}_\ell(q)$ at its rational prime
$\ell\mid a_s$.

\subsection{Reciprocity and the local terms}

Reciprocity expresses $(q/J)_m$ in terms of local symbols above
$q$ and $s$ and at the archimedean places.
In Sharifi's convention, the local norm-residue symbols
$(\beta,q)_{m,v}$ have product $1$ over all places $v$ of $K$.
At a prime ideal $P\mid J$ the local symbol is
$(q/P)_m^{-v_P(\beta)}$. At the other finite places away from
$sq$ it is $1$. Hence
\[
(q/J)_m =
\biggl(\prod_{P\mid q}(\beta,q)_{m,P}\biggr)
\biggl(\prod_{P\mid s}(\beta,q)_{m,P}\biggr)
\biggl(\prod_{v\mid\infty}(\beta,q)_{m,v}\biggr).
\]
This is the product formula used in Sharifi's calculation with
cyclotomic norms \cite[\S3, eq.~(5) and \S\S4,8]{SharifiCycl}.

The archimedean factors are $1$. For $m=2$ this follows from $q>0$,
and for $m>2$ all the archimedean places are complex. The
congruences on $q$ also make the factors above $s$ equal $1$.

\paragraph{Odd $s$.}
Since $s\mid r\mid q-1$ and $s\nmid N$, the gap equations give
$v_s(A)=v_s(q-1)>0$. By \eqref{eq:valuation-coupling},
\[
v_s(q-1)=v_s(2^{\operatorname{ord}_s(2)}-1)+v_s(M)
  \ge b+1.
\]
The $s$-adic logarithm sends $1+s\mathbb Z_s$ to $s\mathbb Z_s$
and sends $q$ into $s^{b+1}\mathbb Z_s$. Dividing $\log q$ by
$s^b$ gives the logarithm of an $m$-th root of $q$ in $\mathbb Q_s$.
Thus the factors above $s$ equal $1$.

\paragraph{$s=2$.}
Since $q\equiv1\pmod8$, we have $r\mid(q-1)/2$ by quadratic
reciprocity, so $q\equiv1\pmod{2m}$. The logarithm on
$1+4\mathbb Z_2$ gives $q=u^{m/2}$ for $u\in1+4\mathbb Z_2$.
Adjoining an $m$-th root of $q$ over a completion of $K$ at $2$
amounts to adjoining $\sqrt u$. Over $\mathbb Q_2$ this gives
the trivial extension for $u\equiv1\pmod8$ and the unramified
quadratic extension for $u\equiv5\pmod8$. The extension remains
unramified after passing to the completion of $K$. Units are
norms in an unramified extension, and $\beta=B-\zeta_m$ is a unit
at $2$ because $B$ is even. Thus the factors above $2$ equal $1$.

The residue-symbol product has now reduced to
\[
(q/J)_m=\prod_{P\mid q}(\beta,q)_{m,P}.
\]
The congruences used at $s$, taken at the largest power $s^b\mid r$
for every prime $s\mid r$, also give $E\operatorname{rad}(r)\mid q-1$.
Here $\operatorname{rad}(r)$ is the product of the distinct primes
dividing $r$.

\subsection{Evaluating the product modulo q}

The primitive $m$-th roots in $\mathbb F_q$ are the $B^i$ with
$1\le i<m$ and $\gcd(i,m)=1$. Use these values of $i$ throughout
this calculation. The prime ideals above $q$ are
\[
P_i=(q,\zeta_m-B^i)
    =q\mathcal O_K+(\zeta_m-B^i)\mathcal O_K.
\]
Indeed, reducing $\mathcal O_K=\mathbb Z[\zeta_m]$ modulo $q$
gives $\mathbb F_q[X]/(\Phi_m(X))$. The distinct linear factors
$X-B^i$ give one copy of $\mathbb F_q$ each, with reduction
$\zeta_m\mapsto B^i$ at $P_i$. Thus $\mathcal O_K/P_i\simeq
\mathbb F_q$ and $(q)=\prod_i P_i$. In particular $v_{P_i}(q)=1$,
so $q$ is a uniformizer at each $P_i$.

\paragraph{The residues of the unit parts.}
At $P_i$, the residue of $\beta$ is $B-B^i$. Among the primes
above $q$, only $P_1$ divides $(\beta)$. Because its ideal norm
is $q$, the norm identity $\operatorname{Norm}_{K/\mathbb Q}(\beta)
=H_s$ gives $v_{P_1}(\beta)=v_q(H_s)=\rho$.

In the completion at $P_i$, dividing $\beta$ by
$q^{v_{P_i}(\beta)}$ leaves a unit. Let $U_i\in\mathbb F_q^\times$
be its residue. For $i\ne1$ the valuation is zero, so $U_i=B-B^i$.
At $P_1$, $U_1$ is the residue of $\beta/q^\rho$.

To find $U_1$, divide the norm identity $H_s=\prod_i(B-\zeta_m^i)$
by $q^\rho$ and reduce at $P_1$. This gives
\[
a_s=U_1\prod_{i\ne1}(B-B^i)\qquad\text{in }\mathbb F_q.
\]
The product on the right is $\Phi_m'(B)$, the formal derivative
of $\Phi_m(X)$ evaluated at $B$. Indeed, over $\mathbb F_q$,
$\Phi_m(X)=(X-B)\prod_{i\ne1}(X-B^i)$, so at $B$ only the term
from differentiating $X-B$ survives. The factors $B-B^i$ are
nonzero, and we obtain
\[
U_1=\frac{a_s}{\Phi_m'(B)}\qquad\text{in }\mathbb F_q.
\]

\paragraph{Multiplying the symbols.}
At $P_i$ the tame formula reduces $(\beta,q)_{m,P_i}$ to
$\bigl((-1)^{v_{P_i}(\beta)}U_i\bigr)^{(q-1)/m}$.
The sign is $1$ since $v_{P_i}(\beta)$ is $0$ or the even $\rho$.
To multiply these symbols in $\mathbb F_q$, use the same reduction
for all of them. A symbol $\zeta_m^k$ gives $B^{ik}$ at $P_i$
and $B^k$ at $P_1$. Thus raising the residue at $P_i$ to the
inverse $u_i$ of $i$ modulo $m$ gives its residue at $P_1$.
Use $u_1=1$. The reduction of $(q/J)_m$ at $P_1$ is therefore
\begin{equation}\label{eq:reciprocity-test}
T_m=
\left(
\frac{a_s}{\Phi'_m(B)}
\prod_{\substack{1<i<m\\ \gcd(i,m)=1}}(B-B^i)^{u_i}
\right)^{(q-1)/m}
\quad\text{in }\mathbb F_q^\times.
\end{equation}
Every solution requires $T_m=1$ for each prime power $m\mid r$.
This test uses $a_s$ only modulo $q$. Compute $\Phi_m(B)$ modulo
$q^{\rho+1}$, then divide by $q^\rho$.

\subsection{The quadratic and cubic tests}

If $2\mid r$, take $m=2$. The product in \eqref{eq:reciprocity-test} is empty
and $\Phi'_2(B)=1$. The condition $T_2=1$ says that $a_2$ is a
square modulo $q$.

If $3\mid r$, take $m=3$. The two residues of $\zeta_3$ above $q$ are $B$ and $B^2$,
with $u_2=2$. Since $B^2+B+1=0$ in $\mathbb F_q$,
\[
B-B^2=2B+1=\Phi'_3(B).
\]
The expression raised to $(q-1)/3$ in \eqref{eq:reciprocity-test}
is therefore $a_3(2B+1)$. In a solution this must be a cube modulo $q$.

This condition also restricts $h=M/E$. Since $a_3=a/(B-1)$ and
$2B+1=B-B^2$, the cubic expression is $-Ba$ modulo $q$.
The two gap equations give
\[
A=a\sum_{j=0}^{h-1}2^{jE}\equiv ah\pmod q,\qquad
A2^{m_2}\equiv-1\pmod q.
\]
Thus $-Ba\equiv B2^{-m_2}h^{-1}\pmod q$.
By \eqref{eq:cubic-return}, $2^{(q-1)/3}\equiv1\pmod q$, so
$2$ is a cube modulo $q$, as are $B$ and $2^{m_2}$.
Consequently, if $3\mid r$, then $h$ must be a cube modulo $q$.

For a numerical check, take $q=3511$, which already fails
the earlier congruences on $q$. It has $r=1755$, $\rho=2$ and
$E=3510$. At $m=3$ the local calculation applies since
$3511\equiv1\pmod9$. Using $B=2^{1170}$ gives
\[
B\equiv2754,\qquad
a_3=\frac{B^2+B+1}{3511^2}\equiv2302\pmod{3511}.
\]
Thus $a_3(2B+1)\equiv-14$ and $T_3=(-14)^{1170}\equiv756\pmod{3511}$.
The test therefore predicts a factor of $a_3$ at which $3511$ has
order divisible by $3$. One such prime is $262657$.
Indeed, $2$ has order $27$ there, so $B$ has order $3$ and
$262657\mid a_3$. Directly, $3511^{87552}\equiv512\pmod{262657}$.

The product can equal $1$ by cancellation. For $q=433$ and $m=8$,
the symbols at $\ell=17,241,38737$ are $\zeta_8^6,\zeta_8^2,1$.
Their product is $1$. The first two symbols each give an order
obstruction. These primes occur in $\Phi_8(2)$, $\Phi_{24}(2)$
and $\Phi_{72}(2)$, respectively. Keeping the relative norms
of $2-\zeta_t$ for $t=8,24,72$ separate recovers the two nontrivial
products. Use compatible roots $\zeta_t^{E/m}=\zeta_m$ to compare
them with the grouped product. In general the roots of
$X^{E/m}-\zeta_m$ that have order $t$ form the conjugates over
$K$ of such a $\zeta_t$. Their factors in $B-\zeta_m$ give the
relative norm from $\mathbb Q(\zeta_t)$ to $K$.

The open step is to prove that some such product is nontrivial
for every remaining $q$, or to control the primes for which all
products equal $1$. At $\ell$, the $s^b$-th power-residue
symbol is nontrivial precisely when
$v_s(\operatorname{ord}_\ell(q))>v_s(\ell-1)-b$.
For example, at $\ell=38737$ in the preceding calculation,
$\operatorname{ord}_{38737}(433)=4842$. Its $2$-valuation is $1$,
and $v_2(38737-1)=4$. The eighth-power symbol therefore equals $1$
although the order is even. Keeping the factors separate resolves
cancellation between them. This example exhibits a second reason
for a product to equal $1$.

\subsection{Restrictions on the multiplier h}
\label{sec:reciprocity-h}

Write $h=M/E$ and $c=m_2$. The cubic calculation is one instance
of a general simplification. Since $B$ is a root of $\Phi_m$
modulo $q$, formal differentiation of
$X^m-1=(X^{m/s}-1)\Phi_m(X)$ gives
\[
(B^{m/s}-1)\Phi'_m(B)=mB^{-1}\qquad\text{in }\mathbb F_q.
\]
Together with $a_s=a/(B^{m/s}-1)$ and
$a\equiv-(2^ch)^{-1}\pmod q$, this expresses the quantity inside
the parentheses in \eqref{eq:reciprocity-test} as
\[
-\frac{B}{m\,2^c h}
\prod_{\substack{1<i<m\\\gcd(i,m)=1}}(B-B^i)^{u_i}.
\]
Thus the grouped test prescribes a power-residue class for $h$
given by a function of the root $B$ and the known $2^c$.

For $m=2$ the expression is $(2^{c+1}h)^{-1}$. Since $2$ is
a square modulo $q$, we obtain that $h$ is a square whenever
$2\mid r$. For $m=4$, $B^2=-1$ and $u_3=3$, so the expression is
$-2/(2^ch)$. Here $4\mid M$, hence $4\mid c$ by
Lemma~\ref{lem:mod-five-split}. Also $-1$ is a fourth power since
$q\equiv1\pmod8$. The quartic test is therefore equivalent to
$2h$ being a fourth power modulo $q$.

\paragraph{Quadratic reciprocity on divisors of $A$.}
Every positive divisor of $A$ is a square modulo $q$.
Indeed, for each prime $\ell\mid A$, the order of $q$ modulo $\ell$
divides the odd integer $N$. Thus $(q/\ell)=1$, and
$(\ell/q)=1$ since $q\equiv1\pmod4$. This also gives $h$ square
modulo $q$ for odd $r$: apply it to $a\mid A$ and use
$2^cah\equiv-1\pmod q$.

In fact $2$ is a fourth power modulo $q$. If $4\nmid r$, then
$r\mid(q-1)/4$ since $16\mid q-1$. Otherwise put $x=2^{r/4}$.
Its residue $i$ has order four. The positive integer $x+1$ divides
$2^r-1$ and is prime to $q$, so $x+1\mid A$. Hence $1+i$ is a
square. Also $i$ is a fourth power since $16\mid q-1$. The identity
$(1+i)^2=2i$ gives the assertion. Consequently
\[
 v_2(r)\le c-2,\qquad 4\mid r\ \Longrightarrow h\in(\mathbb F_q^\times)^4 .
\]
The second assertion uses the preceding quartic test on $2h$.

At $8\mid r$ the grouped test goes one step further. For $m=8$
put $i=B^2$, so $B^4=-1$. The inverse residues of $1,3,5,7$
are themselves modulo eight. Cancelling the derivative in
\eqref{eq:reciprocity-test}, the element tested for an eighth power is
\[
 \frac{a}{B^4-1}(B-B^3)^2(B-B^5)^4(B-B^7)^6
      =-8a(1-i)^8 .
\]
Here $1+i=i(1-i)$. Using $2^cah=-1$ and $4\mid c$, this
condition is equivalent to $2h$ being an eighth power: the quotient
$2^{c-4}$ is an eighth power since $2$ is a square.
Thus if $8\mid r$ and $v_2(r)=c-2$, an eighth-power multiplier
(in particular $h=1$) is excluded. In this case $2$ is a fourth
power and has no eighth root modulo $q$.

\paragraph{An odd conductor and the two factors of its norm.}
Let $m=s^b\mid r$ with $s$ odd, $D=2^{E/(2m)}$, and $B=D^2$.
Choose $C=D$ or $-D$ so its residue $z$ has order $m$ modulo $q$:
$D^2$ has order $m$ and $m$ is odd. Then $z^2=B$. The factorization
\[
 B-\zeta_m^2=(C-\zeta_m)(C+\zeta_m)
\]
takes place over $\mathbb Q(\zeta_m)$. At the prime given by
$\zeta_m\mapsto z$, the first factor vanishes and the second is a unit.
In fact $C+\zeta_m$ is a unit at \emph{every} prime above $q$, since
$-z$ has order $2m$. Its norm $\Phi_m(-C)$ divides
$C^{2m}-1=2^E-1$, hence divides $a$.

Apply the product formula of Section~\ref{sec:reciprocity} to
$C+\zeta_m$. The local symbol at $s$ is trivial since $q$ is an
$m$-th power in $\mathbb Q_s$. Away from $sq$, the primes dividing
the element have residue field $\mathbb F_\ell$ by reduction
$\zeta_m\mapsto-C$. The order of $q$ there divides $N$, prime to
$s$, so their symbols are trivial as well. At the primes above $q$
the unit residues are $z+z^i$. Transporting all symbols to the
embedding $\zeta_m\mapsto z$ gives the necessary condition
\begin{equation}\label{eq:plus-factor-symbol}
 \prod_{i\in(\mathbb Z/m)^\times}(z+z^i)^{u_i}
       \in (\mathbb F_q^\times)^m,\qquad u_i i\equiv1\pmod m .
\end{equation}
This is a separate-factor test of the type above. This particular
factor has no $q$-adic unit part involving $a$ or $h$.

\begin{lemma}\label{lem:fifth-multiplier}
Suppose $5\mid r$. With $F_0=0$, $F_1=1$ and
$F_{k+2}=F_{k+1}+F_k$, one has
\[
       q\mid F_{(q-1)/5},\qquad h\in(\mathbb F_q^\times)^5.
\]
\end{lemma}
\begin{proof}
Use $m=5$ above. Both $2$ and $z$ are fifth powers modulo $q$:
$r\mid(q-1)/5$ by \eqref{eq:ordinary-q-minus-one}, and $25\mid q-1$.
The inverse exponents at $i=1,2,3,4$ are $1,3,2,4$. Thus the
product in \eqref{eq:plus-factor-symbol} equals
\[
 2(1+z)^3(1+z^2)^2(1+z^3)^4
   = -2z(1+z)^2(1+z^2)^5 .
\]
Here $1+z^3=z^3(1+z^2)$ and $(1+z)(1+z^2)=-z^4$.
Consequently $1+z$ is a fifth power, and so is $1+B=1+z^2$.

Put $w=-z^2(1+z)$. Then $w^2=w+1$ and its conjugate root is
$-1/w$. For $t=(q-1)/5$, which is even,
\[
 F_t=0\pmod q\quad\Longleftrightarrow\quad w^{2t}=1
 \quad\Longleftrightarrow\quad w^t=1 .
\]
Indeed the two roots are distinct and $w^t$ has order dividing $5$.
This gives the Fibonacci criterion, equivalent to $1+z$ being a
fifth power.

Finally, write the expression inside the parentheses of
\eqref{eq:reciprocity-test} as
$a_5\prod_{i\ne1}(B-B^i)^{u_i-1}$, where $a_5=a/(B-1)$.
Using $1-B^3=-B^3(1-B^2)$ gives exactly
\[
 \frac{a}{B-1}(B-B^2)^2(B-B^3)(B-B^4)^3
   = a\,\frac{(1-B^2)^5}{1+B}.
\]
Its being a fifth power now says that $a$ is a fifth power. Since
$2^cah=-1$ in $\mathbb F_q$, so is $h$.
\end{proof}

At conductor five the two tests together are equivalent to the two
conditions in Lemma~\ref{lem:fifth-multiplier}. Failure of the
Fibonacci condition excludes every $h$. When it holds, $h=1$
satisfies both tests. The higher congruence at $q$ below is a separate condition.

\paragraph{A lower gap bound when $h=1$ and $5\mid r$.}
In this case $q\equiv1\pmod5$ forces $4\mid M$. Since
$M=E=\operatorname{lcm}(2,r)$, we have $4\mid r$, $M=r=20t$.
Put $x=2^t$. Its residue has order $20$ modulo $q$, so the full
$q^\rho$ in $2^r-1$ divides
\[
 \Phi_{20}(x)=x^8-x^6+x^4-x^2+1<x^8.
\]
It follows from $q^{N+\rho}>2^{r+c}$ that $N>3\rho/2$.

For $t$ odd this improves to $N>4\rho$. The case $t=1$ is
excluded by $q^\rho\le\Phi_{20}(2)=205$ and $q\ge97$.
Otherwise $x\ge8$ and $z=2^{(t+1)/2}$ satisfies $z^2=2x$.
For $\epsilon=\pm1$ define
\[
 P_\epsilon=x^4+\epsilon x^3z+x^3-x^2+x+\epsilon z+1.
\]
Multiplication gives
$(x+1-\epsilon z)P_\epsilon=x^5+1+\epsilon x^2z$
and $P_+P_-=\Phi_{20}(x)$. No odd prime divides both factors:
it would divide the difference $2x^2z$ of the two numerators.
Thus $q^\rho\le P_+<x^4(1+\sqrt{2/x}+1/x)\le(13/8)x^4$.
Since $(13/8)^5<16\le2^c$, we get $q^{5\rho}<2^{r+c}<q^{N+\rho}$.

\paragraph{Recovering $h$ from $q$.}
Put $Q=q^\rho$. Since $2^E\equiv1\pmod Q$, the geometric sum for
$A/a$ is $h$ modulo $Q$. Lemma~\ref{lem:sharp-h} gives $N>\rho$,
so the gap equation yields
\[
h\equiv-(2^ca)^{-1}\pmod Q,\qquad 1\le h<q^2\le Q.
\]
There is at most one possible $h$ for each $q$. Computing
$2^E\bmod Q^2$ suffices to recover $a\bmod Q$ and hence this
candidate. The independent bound $h<q^2$ in Lemma~\ref{lem:sharp-h}
is what singles out one integer in the residue class. The bound
\eqref{eq:number-of-returns} and its sharpening instead give
lower bounds on $N$ for this $h$. If $h\ge2$, then
$N>(2h-1)\rho\ge3\rho$, so the next
necessary congruence is
\[
1+2^ca\left(h+Q a\binom h2\right)\equiv0\pmod{Q^2}.
\]
It follows by writing $2^E=1+Qa$ in the geometric sum and retaining
the constant and linear terms in $Q$. If $h=1$, the corresponding
test is $1+2^ca\equiv0\pmod{q^{\rho+1}}$.

\paragraph{What lifting the fifth-power test can detect.}
For $q\ne5$ and $k\ge1$, reduction gives an isomorphism
\[
 (\mathbb Z/q^k)^\times/((\mathbb Z/q^k)^\times)^5
 \ \simeq\ \mathbb F_q^\times/(\mathbb F_q^\times)^5 .
\]
Indeed, the kernel of reduction on units is a $q$-group, on which
taking fifth powers is an automorphism. Equivalently, every fifth
root of a nonzero residue lifts uniquely by Hensel's lemma.
When $h=1$, knowing $1+2^ca\equiv0\pmod{q^\rho}$ fixes $a$
modulo $q^\rho$. Its $q$ lifts modulo $q^{\rho+1}$ have identical
fifth-power class, while exactly one gives
$1+2^ca\equiv0\pmod{q^{\rho+1}}$. Thus lifting those same
fifth-power characters cannot distinguish the required lift.
The integer $a=(2^E-1)/q^\rho$ is fixed; the question is which
lift that specific integer realizes.

More exactly the original equality gives
$2^{Eh}-(1-q^\rho/2^c)=q^{N+\rho}/2^c$. Both units on the
left belong to $1+q^\rho\mathbb Z_q$. Write $\ell_q$ for the
$q$-adic logarithm. It preserves valuations of differences there,
and $v_q(\ell_q(2^E))=\rho$. Consequently
\[
 v_q\!\left(h-
 \frac{\ell_q(1-q^\rho/2^c)}{\ell_q(2^E)}\right)=N.
\]
This identifies the precision required of the actual multiplier.
It uses the full gap equation; an independent estimate for this
valuation could be compared with the real gap bounds.

\paragraph{Recovering $N$ at the prime $2$.}
Put $u=(q-1)/2^c$. Once $h$ is known, $M=hE$ is known, and the
two gap equations give
\[
q^N\equiv 1-2^c q^{-\rho}\pmod{2^{M+c}}.
\]
Already modulo $2^{2c}$, expansion of $q^N=(1+2^cu)^N$ gives
$Nu\equiv-1\pmod{2^c}$. In particular, $Nu+1\ge2^c$.
Combining this with $N\log q<M\log2$ gives
\[
\frac{2^c-1}{u}\log q<hE\log2.
\]
For fixed odd $u$ this now gives $2^c\le2519u+1$ directly.

More generally take $j=\lceil\log_2 M\rceil$. The order of $q$
modulo $2^{c+j}$ is $2^j$ by lifting, and the subgroup generated
by $q$ is the full subgroup congruent to $1$ modulo $2^c$.
Thus the displayed congruence determines $N$ modulo $2^j$ by
lifting one bit at a time. Since $0<N<M\le2^j$, there is at most
one candidate $N$. It must satisfy both the height window and
the congruence at $q$ that follows from the gap equation.

For the congruence at $q$, Lemma~\ref{lem:uniform-N} reduces the
required lower bound to a constant: if the least positive residue
of $-(2^ca)^{-1}$ modulo $q^\rho$ exceeds $630$, the case is excluded.

\paragraph{A size bound from successive lifts at $2$.}
Put $T=2^c$ and $n=N+\rho$. The same gap equation reads
$q^\rho(1-q^N)=T-2^{m_1}$. The first lift gave $2^c\le Nu+1$.
The next one gives
\[
 T^3\mid N^2(q-1)+NT+\tfrac12(n+\rho-1)T^2.
\]
Every term on the right is positive. The same sign persists at all
orders, giving the following uniform inequality.
\begin{lemma}\label{lem:lift-size}
For $c\ge28$, put $k=\lfloor(c-2)/26\rfloor$. Then
\[
 \log_2 q>c+(c-23)k-1.
\]
\end{lemma}
\begin{proof}
Write the formal solution of $X^\rho(1-X^N)=t$ at $X(0)=1$
as $X(t)=1-\sum_{i\ge1}b_i t^i$. To compute the coefficients,
put $W=1-X^N$. Then $W=t(1-W)^{-\rho/N}$ and
$1-X=1-(1-W)^{1/N}$. Lagrange inversion gives
\[
 b_i=\frac{\prod_{j=1}^{i-1}(i\rho-1+jN)}{i!\,N^i}>0 .
\]
In particular $b_1=1/N$, $b_2=(n+\rho-1)/(2N^2)$,
giving the displayed congruence since $m_1\ge3c$.
Also $N^{2i-1}b_i$ is an integer: in the equation for $1-X$
the linear coefficient is $N$ and all coefficients are integers.
For $i>1$ each nonlinear term contributing to $Nb_i$ is a
product of $j\ge2$ earlier coefficients with indices summing to $i$.
Induction bounds its denominator by $N^{2i-j}$.
Since $\rho<N$, the product formula and $i!\ge(i/\exp(1))^i$
give $b_i\le6^i$.

For $1\le l\le k$ set $D=N^{2l-1}$. Whenever $m_1\ge(l+1)c$,
substitution at $t=T-2^{m_1}$ gives
\[
 T^{l+1}\mid D(q-1)+D\sum_{i=1}^l b_iT^i.
\]
Indeed the series converges at this even $t$ in $\mathbb Z_2$;
its value is $q$ by uniqueness of the odd root (the derivative
at $1$ is $-N$). Now $6N^2<2^{26}$ and $4(6N^2)^l\le T$.
The sum including its factor $D$ is at most
$2(6N^2)^lT^l\le T^{l+1}/2$. Positivity and divisibility yield
$q>T^{l+1}/(2N^{2l-1})$.

The precision hypothesis holds successively. Initially
$2^{m_1}=q^n-q^\rho+T>q^{n-1}>T^{n-1}$, with $n\ge7$.
After step $l$ we have $q>T^l$ since $2N^{2l-1}<T$.
Thus $2^{m_1}>T^{l(n-1)}\ge T^{l+2}$ for the next step.
Take $l=k$ and use $N^2<2^{23}$.
\end{proof}

\paragraph{An effective bound for a prescribed range of $c$.}
For fixed $N,\rho,c$ put $P(X)=X^{N+\rho}-X^\rho+2^c$.
A multiple root $z$ would be nonzero and satisfy both
$z^N=\rho/(N+\rho)$ and $2^c=Nz^\rho/(N+\rho)$, impossible
since $|z|<1$. Schinzel and Tijdeman's theorem
\cite{SchinzelTijdeman} therefore gives an effectively computable bound
on $m_1$ in $P(q)=2^{m_1}$. Since $q^{n-1}<2^{m_1}$ and
there are finitely many $(N,\rho)$, this bounds $q$ effectively
for every prescribed range $c\le c_0$. The bound depends on $c_0$.
Lemma~\ref{lem:lift-size} gives $c=O(\sqrt{\log q})$ for the
complement. The explicit Schinzel--Tijdeman estimate of
Bérczes, Evertse and Győry \cite[Thm.~2.3]{BEG} specializes to
$m_1\le (10n^2)^{40n}2^{11nc}$: the ground field is $\mathbb Q$,
the only place in $S$ is real, and the coefficient height is $c\log2$.
This grows too fast to meet Lemma~\ref{lem:lift-size}.

\subsection{The odd part of the binary gap}
\begin{lemma}\label{lem:odd-gap-part}
In the remaining case the odd part of $M$ is neither $1$ nor $3$.
\end{lemma}
\begin{proof}
Write $M=2^a b$ with $b=1$ or $3$, and set $D=2^{256b}-1$.
The only possible repeated prime factor of $D$ is $3$.
For $b=1$ this follows from the factorizations of the Fermat numbers
$F_i=2^{2^i}+1$, $0\le i\le7$ \cite{BrentFermat}.
For $b=3$ the additional factors are those of
$\Phi_{3\cdot2^i}(2)$, $0\le i\le8$ \cite{CunTables}.
The certificate below verifies both products and the primality of every factor.

Since $q^\rho\mid2^M-1$, $\rho\ge2$ and $q\nmid M$, the order
$r=\operatorname{ord}_q(2)$ cannot divide $256b$: otherwise lifting
at $r$ gives $q^\rho\mid D$. Thus $512\mid r$ and $256b\mid M$.
Any prime $\ell$ with $512\mid\operatorname{ord}_\ell(2)$
is $1$ modulo $1024$. Indeed $\ell\equiv1\pmod{512}$, so $2$ is
a quadratic residue modulo $\ell$ and its order divides $(\ell-1)/2$.
In particular $c\ge10$. Every prime factor of $2^M-1$ outside $D$
has order divisible by $512$, so it contributes $1$ to $u\bmod1024$,
where $u=(q-1)/2^c$. We have $Nu\equiv-1\pmod{1024}$.

For $\ell^e\Vert D$, lifting across the power of two $M/(256b)$
gives $v_\ell(A)=e$, since $q\nmid D$.
If $\gcd(N,\ell-1)=1$, then $q^N\equiv1\pmod\ell$ forces $\ell\mid u$.
Whenever $\ell\mid u$, its exponent in $u$ is exactly $e-v_\ell(N)$
and must be positive. Otherwise its contribution is $1$.

For $b=1$ the calculation can be done by hand. All factors of
$F_0\cdots F_5$ must occur in $u$ exactly once: for $F_0,\ldots,F_4$
the predecessors are powers of two; $F_5=641\cdot6700417$ and the
odd parts of their predecessors are $5$ and $3\cdot17449$, respectively.
Here $17449$ is prime, $N\le2519$, and $\gcd(N,30)=1$ since $4\mid M$.
The factors of $F_i$ for $i\ge6$ are $1$ modulo $256$ by the same
quadratic-residue argument. Hence $u\equiv2^{64}-1\equiv-1\pmod{256}$.
So $N\equiv1\pmod{256}$, leaving $769,1537,1793$ (the value
$257$ is excluded by $v_{257}(u)=1$ and $e=1$).
For each of these three integers, each prime factor $\ell$ of
$F_6F_7$ has $\gcd(N,\ell-1)=1$. Thus all occur in $u$ and
$u\equiv2^{256}-1\equiv-1\pmod{1024}$, contrary to all three values of $N$.

For both $b$ the same finite calculation starts with the residue $1$,
multiplies by the permitted contributions of each $\ell^e\Vert D$,
and retains $N$ if $-N^{-1}$ occurs. The full factorization of $D$
is used in each column, including primes whose presence rules out
$N$ by valuation alone.
\[
\begin{array}{c|c|c|c}
b & \bmod 256 & \bmod 512 & \bmod 1024\\ \hline
1&769,1537,1793&1537&\varnothing\\
3&1537,1793&1537,1793&\varnothing
\end{array}
\]
The PARI/GP certificate is \path{scripts/check_odd_part.gp}.
\end{proof}

For general odd $b$, $2^{256b}-1$ still contains every prime of
$2^M-1$ that can affect $u\bmod1024$ when $256\mid M$.
If the order of $q$ divides $256b$, first retain the possible
$q>3$ whose square divides $2^{256b}-1$. Otherwise $c\ge10$
and the same residue calculation applies. For varying $b$,
the missing step is to control the factors of this smaller Mersenne
number uniformly.

The necessary factor can sometimes be located at a smaller exponent.
Put $K=(q^N-1)/(q-1)$ and let $H$ be the product of the full
$\ell$-powers in $N$ for primes $\ell\mid q-1$. Lifting shows that
$K/H$ has only primes at which $q$ has order greater than $1$.
Also $K/H\equiv N/H\pmod{2^c}$. If $N/H\not\equiv1\pmod{2^s}$
for some $s\le c$, at least one $\ell\mid K/H$ differs from $1$
modulo $2^s$. For $s\ge4$, quadratic residuosity as above gives
$v_2(\operatorname{ord}_\ell(2))\le s-2$. Thus already
\[
 \ell\mid 2^{\,b\,2^{\min(a,s-2)}}-1,\qquad
 1<\operatorname{ord}_\ell(q)\mid\gcd(N,\ell-1).
\]
Here $\ell\mid K\mid A$ supplies $\operatorname{ord}_\ell(2)\mid M$.
This identifies which factors could prevent the product argument
from excluding another odd part $b$.

In fact $H=1$ except possibly when $N=1093$ and $1093\mid q-1$.
For a shared odd prime $\ell$, we have $\ell\nmid M$ by
$\gcd(g,q-1)=1$. Then $v_\ell(2^M-1)=v_\ell(q-1)+v_\ell(N)\ge2$,
so $2^{\ell-1}\equiv1\pmod{\ell^2}$. Checking this congruence at
the odd primes up to $2519$ gives just $1093$ (same certificate);
its only odd multiple in the range is itself. In particular, if
$4\mid M$, $c\ge10$, and $H=1$, take
$s=\max(4,1+v_2(N-1))$. This is at most $10$: the only integers
in $[5,2519]$ congruent to $1$ modulo $1024$ are $1025$ and
$2049$, both excluded by $\gcd(N,30)=1$.

\paragraph{Using the same multiplier for several orders.}
Let $t\mid N$ be prime, $t\ne q$, $t\nmid q-1$, and
$\gcd(t-1,N)=1$. Then $\operatorname{ord}_t(2)\nmid M$:
otherwise $t\mid A$, so the order of $q$ modulo $t$ divides
$\gcd(N,t-1)=1$. The gcd hypothesis holds in particular for
the least prime divisor of $N$.

More generally, suppose $d\nmid M$. For every odd prime $P$
with $d\mid\operatorname{ord}_P(2)$, we have $P\nmid q-1$:
if $P\mid q-1$, then $P\mid u\mid A$ and
$\operatorname{ord}_P(2)\mid M$. Since $E\mid q-1$, the full
power of every such $P$ in $M=hE$ must divide $h$. Their product
therefore divides $h$.

Take $N=1793=11\cdot163$. Neither prime divides $q-1$ by the
preceding valuation check, and neither can be $q$ since
$q\equiv1\pmod{48}$. Their predecessors are coprime to $N$,
so $10=\operatorname{ord}_{11}(2)\nmid M$ and
$162=\operatorname{ord}_{163}(2)\nmid M$. As $M$ is even,
$5\nmid M$ and $81\nmid M$. In particular $331$ cannot divide
$A$, since $\operatorname{ord}_{331}(2)=30$.

For each $t=11,163$, the factor $\Phi_t(q)\equiv t\equiv3\pmod4$
requires a prime $\ell\equiv3\pmod4$ with $\operatorname{ord}_\ell(q)=t$.
Consider assigning $23$ to $11$ and $13367$ to $163$. Both are
primes in the required classes ($13367-1=82\cdot163$).
Their orders at base $2$ are $11$ and $41$, forcing $11\cdot41\mid M$.
Since $5$ divides both $\operatorname{ord}_{11}(2)=10$ and
$\operatorname{ord}_{41}(2)=20$, the preceding argument gives
$451\mid h$, whereas $h<(1793+2)/4$. Thus $451\nmid M$ for this
$N$, and in particular this assignment is impossible.

\paragraph{How much of $q-1$ can come from primes dividing $M$?}
Let $V$ be the largest divisor of $u$ coprime to $M$ and $L=\log_2 q$.
Then
\[
 L<1+c+\log_2 V+\sqrt M+2\log_2 M,\qquad M<5037L.
\]
Indeed, for each prime $\ell\mid\gcd(u,M)$ put $j=v_\ell(M)$ and
$d=\gcd(\ell-1,M/\ell^j)$. Its order at base $2$ divides the even
integer $d$. Lifting gives
$v_\ell(u)+v_\ell(N)=v_\ell(2^d-1)+j$. The full $\ell$-power
in $2^d-1$ belongs to one of $2^{d/2}\pm1$, so its binary
logarithm is at most $1+d/2$. Put $z=\sqrt M$. For $3\le\ell\le z$
use that $x/\log x$ increases on $[3,\infty)$; for $\ell>z$
use $M/\ell\le z$. In both cases $d\le z\log\ell/\log z$.
Summing over $\ell$ gives $\sum d\le2\sqrt M$, and both
$\sum j\log_2\ell$ and the number of terms are at most $\log_2 M$.
This proves the bound since $q=1+2^c u$.

For example, $V\le2^c$ implies $c<448$ and $L<7168$.
For $L\ge7168$ one has
$1+\sqrt{5037L}+2\log_2(5037L)<0.85L$ (divide by $L$;
each term decreases). Hence $L<(40/3)c$ in this range.
If $c\ge448$, Lemma~\ref{lem:lift-size} gives $L>17c>7168$,
a contradiction. With $c<448$, $L\ge7168$ is then impossible as well.
In general the same inequality gives $\log_2 V\ge L-O(\sqrt L)$
using Lemma~\ref{lem:lift-size}. There is no upper bound on $V$
in terms of $c$ here.

\subsection{Computational supplements}

\begin{itemize}
\item \path{md/mersenne-tables.md} gives orders, the factors used in $D_E$,
and Frobenius in the cyclotomic fields for illustrative $q$.
The GAP source is \path{scripts/mersenne_tables.g}. For the cyclotomic
Frobenius formula see Milne \cite{MilneANT}.
\item \path{md/order-obstructions.md} checks all prime choices with
$E\le256$ using \path{scripts/check_order_obstruction.g} and
\path{scripts/mersenne-factors-256.g}. There are 395 certificates among
396 choices, with $q=3$, $a=1$ the exception. All have $\rho=1$.
\item \path{md/cyclotomic-reciprocity.md} gives the derivation and examples.
\path{md/cyclotomic-reciprocity-tables.md} compares the calculation at $q$
with direct symbols at the prime factors and includes examples
with $\rho=2$. The PARI/GP functions are in
\path{scripts/cyclotomic_reciprocity.gp}. The tables come from
\path{scripts/check_cyclotomic_reciprocity.gp}.
\item \path{scripts/inspect_reciprocity.sage} compares the grouped
product with its separate cyclotomic factors and optionally checks
the individual prime orders. It distinguishes cancellation from
an order that the symbol misses. Its function \texttt{gap\_candidate}
computes the candidate $h$ and checks the next gap congruence.
\end{itemize}

\section{Exercises and remaining questions}
\label{sec:exercises}
% Consolidated development plan. Closed calculations and open targets are
% identified within each item; no blanket coprimality claim is assumed.
\begin{enumerate}
\item \textbf{Cyclotomic multiplicities.}
With $r=\operatorname{ord}_q(2)$ and $\rho=v_q(2^r-1)$, prove
\[
v_q(\Phi_d(2))=\begin{cases}
\rho&d=r,\\1&d=rq^j,\ j\geq1,\\0&\text{otherwise}.
\end{cases}
\]
Remove these powers to obtain $A$ as a product of cyclotomic factors.
Then remove each intrinsic factor to recover the pairwise coprime
primitive parts in \cite[\S2]{PomCycl}. Determine which Aurifeuilian
factor must enter $A$ when $q$ divides the other.

\item \textbf{Orders including multiplicity.}
For $t=\operatorname{ord}_\ell(q)$, prove
\[
\operatorname{ord}_{\ell^e}(q)
=t\ell^{\max\{0,e-v_\ell(q^t-1)\}}.
\]
Using Section~\ref{sec:exact-orders}, recover $M$ and $N$ as the lcms
of the orders modulo the prime-power factors of $A$.
Work out the restrictions supplied by primitive Mersenne factors at
$2,3,4,8,12$, including the exception when the factor is $q$.
Use $\gcd(N,6)=1$ to sharpen the residue classes. Starting with
Section~\ref{sec:minimum-returns}, propagate selected prime-power
orders in both directions. Seek families where the forced residue
product exceeds $(q-1)/2^{m_2}$, or the minimum return exceeds the
height bound. Distinguish a fixed point of these restrictions from
an actual solution of the gap equations.

\item \textbf{Coprime gaps: multipliers.}
Assume $\gcd(M,N)=1$. For primitive factors $\ell=jM+1\neq q$,
use $\operatorname{ord}_\ell(q)\mid\gcd(j,N)$ to exclude specified
families of multipliers. Determine which odd prime factors of $N$
are forced by several such $j$ when $\ell\nmid q-1$.
Form $D$ from Section~\ref{sec:forcing-common-factor}; compare its
prime powers with the odd part of $q-1$. Read the index condition
as avoidance of splitting in $\mathbb Q(\zeta_s,2^{1/s})$, $s\mid N$.
Open target: force $D\nmid(q-1)/2^{m_2}$ for a specified family of $M,N$.

\item \textbf{Common gaps: full cyclotomic factorization.}
For $d\le503$ with at most two prime factors, both at least $11$,
compare the exact bound from $H_d\mid A$ with
Section~\ref{sec:full-gcd}, retaining $2^{m_2/2}q^{n_1/2}$ in the
denominator. Identify the prime-power common divisors untouched by
this size argument.
Use Section~\ref{sec:additional-returns}, where $dE\mid M$ is now
necessary. Use $(2h-1)n_1<N$ and extend the residue table beyond
$N=121$. For $N=5d$, compute the finite-field polynomial gcd and
look for a condition on $d$ making it $1$ when $d\equiv1\pmod3$.
Use \eqref{eq:ordinary-q-minus-one} to impose
the additional restrictions from odd prime factors of $q-1$.

\item \textbf{Descent and carries.}
For fixed odd prime $d$, classify the divisors of $\tau^d+1$ which
produce integer $U,V$ in \eqref{eq:quotient-divisor}. Verify the example
in Section~\ref{sec:descent} and its failure of the full gap equations.
Complete the $2$-adic cases omitted by $a>m_2+v_2(d-1)$, including
possible cancellation at equal valuations. Open target: a decreasing
transformation preserving the prime-power bases and full gap equations,
or a direct exclusion from the digit and odd-part conditions.

\item \textbf{Factors remaining in C.}
For a forced $\ell^e\mid A$, determine
$\ell^{\max\{0,e-v_\ell(Y-1)\}}\mid C$.
Use the mixed-base gcd formula to combine coprime contributions.
Open target: a lower bound exceeding $Y^{s/2}/[2q(Y-1)]$ for a
specified family. Test $M=2s^a$ and $M=4s^a$ with fixed $a\geq2$,
imposing $s\mid N$ when using $C$.

\item \textbf{The remaining valuation branch.}
Check the analytic bounds and the integer certificate in
Section~\ref{sec:q-divides-exclusion}. For the remaining case,
$q\nmid M$ and $\rho=n_1$ is even. Use the other cyclotomic factors
to restrict multiples of $r$, retaining
$2^{(q-1)/3}\equiv1\pmod{q^2}$. An effective absolute bound on $q$ remains open.
\end{enumerate}


\section{An effective cutoff for the prime base}
\label{sec:effective-q-cutoff}

Use the branch $A>1$, $q>3$, $q\nmid M$. Write $c=m_2$, $e=n_1=\rho$,
$n=n_2=e+N$, $T=2^c$, and $Q=q^e$. Here $e\ge2$ is even, $n$ is odd,
$e<N\le2519$, and $T\le q-1$. With $M=hE$, take the square root
$X=2^{hE/2}$ in the gap equation:
\[
X^2=1-\frac{Q}{T}+\frac{q^n}{T}.
\]
This is the equation underlying the congruence that recovers $h$
in Section~\ref{sec:reciprocity-h}. The square root works also for $h=1$.

Set $k=\lfloor n/(2e)\rfloor$, so $1\le k\le630$ and
\[
2ek\le n-1,\qquad (n+1)/2\le (k+1)e\le n.
\]
The truncated square root
\[
S=\sum_{i=0}^{k}\binom{1/2}{i}(-Q/T)^i
\]
satisfies $S^2\equiv1-Q/T\equiv X^2\pmod{Q^{k+1}}$. The
denominators are powers of two. Since $S\equiv1\pmod q$, exactly
one of $X-S$ and $X+S$ is divisible by $q$. For the corresponding
sign $\varepsilon$, $X-\varepsilon S\equiv0\pmod{Q^{k+1}}$.

Put $B=Q/T>q$. The binomial coefficients have absolute value at
most one, so $|S|<2B^k$. Also $X^2>q^n/(2T)$, and hence
\[
\frac{X^2}{B^{2k}}
>\frac{q^{n-2ek}T^{2k-1}}2>4.
\]
Thus $X>|S|$. The integer $L=(4T)^k$ clears the denominators of $S$,
since $4^i\binom{1/2}{i}$ is integral. The nonzero integer
$L(X-\varepsilon S)$ is divisible by $Q^{k+1}$. Using $X<q^{n/2}$
gives
\[
q^{(k+1)e}\le |L(X-\varepsilon S)|<2(4T)^kq^{n/2},
\qquad
\log_2 q<2+2k(c+2)\le1260c+2522.
\]

Lemma~\ref{lem:lift-size} gives
\[
\log_2 q>c+(c-23)\left\lfloor\frac{c-2}{26}\right\rfloor-1
\quad(c\ge28).
\]
If $c\ge32768$, the floor is at least $1260$ and the lower bound
is at least $1261c-28981>1260c+2522$. Therefore $c<32768$, giving
\[
\boxed{q<2^{41\,290\,202}.}
\]
Section~\ref{sec:q-divides-exclusion} treats the other branch of the
effective argument.

\section{Quadratic norms modulo two}
\label{sec:quadratic-norm}

\begin{lemma}\label{lem:quadratic-norm}
For primes $p,q$, two distinct partitions with positive exponents
force $(p,q)=(11,3)$.
\end{lemma}
\begin{proof}
Lemmas~\ref{lem:Aone} and~\ref{lem:qthree} treat $A=1$ and $q=3$.
Suppose $q>3$. Write $e=n_1$ and $n=n_2$. By
Lemma~\ref{lem:new-parities}, $m_1,m_2,e$ are even and $n$ is odd. Set
\[
x=2^{m_1/2},\qquad y=2^{m_2/2},\qquad D=q^e-2^{m_2}.
\]
Then $D$ is odd and $D>1$, since $2^{m_2}\le q-1$ by
Lemma~\ref{lem:valuations} and $e\ge2$. The two partitions give
\[
x^2+D=q^n,\qquad y^2+D=q^e.
\]

Work in $R=\mathbb Z[w]$ with $w=\sqrt{-D}$. Choose the sign of $y$
so that $q\nmid x+y$. Such a sign exists since $q\nmid2x$. Take
$\alpha=x+w$ and $\beta=y+w$, with respective norms $q^n$ and $q^e$.
Put $H=q^{ne}$ and $U=\alpha^e\overline\beta^{\,n}$.
In $R/HR$, the sum $\overline\alpha+\beta=x+y$ is a unit by
Bézout over $\mathbb Z$. Thus $\overline\alpha^{\,e}$ and
$\beta^n$ generate the unit ideal in $R/HR$. Both annihilate $U$
there, because
\[
\alpha^e\overline\alpha^{\,e}=\beta^n\overline\beta^{\,n}=H.
\]
It follows that $U=Hz$ for some $z\in R$. Taking norms gives
$\operatorname{Norm}(z)=1$. Since $D>1$, the equation
$a^2+Db^2=1$ for $z=a+bw$ forces $z=\pm1$.

Modulo $2R$, both $\alpha$ and $\overline\beta$ reduce to $w$.
We have $w^2=1$ in this quotient, so $U$ reduces to $w^{e+n}=w$.
But $Hz$ reduces to $1$ because $H$ is odd and $z=\pm1$.
The basis $1,w$ of $R$ gives $w\ne1$ modulo $2R$, a contradiction.
The pair $(11,3)$ is attained by $11=2^3+3=2+3^2$.
\end{proof}

This is the quadratic-order argument formalized in \path{Quadratic.lean}.
It works with $D$ itself, even when $D$ has square factors.

\section{Reading the proof}
\label{sec:proof-reading}

\begin{enumerate}
\item Section~\ref{sec:setup} gives the common factor $A$, the
identity for $Ap$, the gap parities, and the case $A=1$.
\item From Section~\ref{sec:bounds}, use $m_2=v_2(q-1)$ from
Lemma~\ref{lem:valuations} and the half-gap bound
$q^{n_1}\le2^{M/2}+1$ from Lemma~\ref{lem:bounds}.
Also use $q\equiv1\pmod3$ for $q>3$: since $M$ is even,
$3\mid A$ when $q\ne3$, so $q^N\equiv1\pmod3$. The gap $N$ is odd.
\item In Section~\ref{sec:parities}, read the proof of the base
$3$ case and the parity argument for $m_1,m_2,n_1,n_2$.
\item Section~\ref{sec:quadratic-norm} closes the $q>3$ case.
\end{enumerate}

The rest of the size and residue estimates in
Sections~\ref{sec:bounds} and~\ref{sec:parities}, and the route
from Sections~\ref{sec:mersenne} through~\ref{sec:exercises}, are
unused in the classification. This includes the common-divisor
analysis, Laurent's bound, the finite check for $q\mid M$, and the
order and reciprocity tests. The effective cutoff in
Section~\ref{sec:effective-q-cutoff} and the appendices are also
independent of the classification proof.

\appendix
\clearpage
\section{Congruence and case reference}
\label{sec:reference-tables}
Assume $A>1$, $q>3$ and the two partitions of the prime $p$.
The proved reductions give $q\nmid M$. Write
\[
r=\operatorname{ord}_q(2),\quad k=M/2,\quad
c=m_2=v_2(q-1),\quad g=\gcd(M,N),\quad E=\operatorname{lcm}(2,r).
\]

\subsection{Residues of q and its square}
{\small\renewcommand{\arraystretch}{1.15}
\begin{longtable}{@{}p{0.49\textwidth}p{0.22\textwidth}p{0.23\textwidth}@{}}
\textbf{Modulus and condition}&\textbf{$q$}&\textbf{$q^2$}\\\hline
\endhead
Modulo $3$&$1$&$1$\\
Modulo $8$&$1$&$1$\\
Modulo $16$&$1$&$1$\\
Modulo $48$&$1$&$1$\\
Modulo $240$&$1,97,193$&$1,49$\\
Modulo $6r$&$1$&$1$\\\hline
$3\nmid M$: modulo $9$&$4,7$&$4,7$\\
$3\mid M$: modulo $9$&$1$&$1$\\
$3\mid M$: modulo $7$&$1$&$1$\\
$4\mid M$: modulo $5$&$1$&$1$\\
$M\equiv2\pmod4$: modulo $5$&$2,3$&$4$\\
$8\mid M$: modulo $17$&$1$&$1$\\
$12\mid M$: modulo $13$&$1$&$1$\\
$16\mid M$: modulo $257$&$1$&$1$\\
$18\mid M$: modulo $19$&$1$&$1$\\\hline
$4\mid M$: modulo $240$&$1$&$1$\\
$M\equiv2\pmod4$: modulo $240$&$97,193$&$49$\\
\hline
\end{longtable}}
The old $q\equiv5\pmod8$ class is excluded by
Lemma~\ref{lem:q-one-mod-eight}. The rows modulo $7$ and $13$ use
$3\nmid N$ and are stronger than the earlier quadratic-residue lists.
Also $q^2\equiv1\pmod{96}$.

\subsection{Valuations}
{\small\renewcommand{\arraystretch}{1.2}\noindent
\begin{tabular}{@{}p{0.43\textwidth}p{0.25\textwidth}p{0.26\textwidth}@{}}
\textbf{Prime}&\textbf{Valuation of $q-1$}&\textbf{Valuation of $q^2-1$}\\\hline
$2$&$c\text{ even},\ c\ge4$&$c+1$\\
$3$&$1+v_3(M)$&$1+v_3(M)$\\
Odd $\ell\mid q-1$, $\eta_\ell=1$&$1+v_\ell(M)$&$1+v_\ell(M)$\\\hline
\end{tabular}\par}
\medskip\noindent
Here $\eta_\ell=v_\ell(2^{\operatorname{ord}_\ell(2)}-1)$.
The last row also forces $\ell\nmid N$; it applies to
$\ell=3,5,7,11,13$ whenever they divide $q-1$.
The exact valuation at the base $q$ is
\[
v_q(2^M-1)=v_q(2^r-1)=n_1\text{ even},\qquad n_1\ge2.
\]

\clearpage
\subsection{Powers of two modulo q and its square}
{\small\renewcommand{\arraystretch}{1.2}
\begin{longtable}{@{}p{0.49\textwidth}p{0.22\textwidth}p{0.23\textwidth}@{}}
\textbf{Expression}&\textbf{Modulo $q$}&\textbf{Modulo $q^2$}\\\hline
\endhead
$2^r$&$1$&$1$\\
$2^M$&$1$&$1$\\
$2^{q-1}$&$1$&$1$\\
$2^{(q-1)/2}$&$1$&$1$\\
$2^{(q-1)/3}$&$1$&$1$\\
$2^{(q-1)/12}$&$1$&$1$\\
$2^{(q-1)/6}$&$1$&$1$\\
$2^{(q-1)/\ell}$, $\ell\mid q-1$ odd and $\eta_\ell=1$&$1$&$1$\\
$2^k$, if $r\mid k$&$1$&$1$\\
$2^k$, if $r\nmid k$&$-1$&$-1$\\\hline
\end{longtable}}
By quadratic reciprocity on divisors of $A$, $r\mid(q-1)/4$.
Together with $r\mid(q-1)/3$, this gives $r\mid(q-1)/12$.
The order remains $r$ modulo $q^2$, so the signs in the last two
rows agree. Also $v_2(r)\le c-2$.

For each prime power $\ell^f\mid A$,
\[
q^N\equiv(q^2)^N\equiv1\pmod{\ell^f},\qquad
\operatorname{ord}_{\ell^f}(q)
=\operatorname{ord}_{\ell^f}(q^2)\mid N.
\]
The equality of orders uses that $N$ is odd. For the full modulus,
\[
\operatorname{ord}_A(2)=M,\qquad
\operatorname{ord}_A(q)=\operatorname{ord}_A(q^2)=N.
\]

\subsection{The prime p and the gap conditions}
{\small\renewcommand{\arraystretch}{1.15}
\begin{longtable}{@{}p{0.35\textwidth}p{0.60\textwidth}@{}}
\textbf{Condition}&\textbf{Restriction}\\\hline
\endhead
$p$ modulo $96$&$65$\\
$p$ modulo $5$&$2,3$; if $4\mid M$, the residue is $2$\\
$p$ modulo $480$&$257,353$\\
$p$ modulo $g$&$1$\\\hline
Exponents&$m_1,m_2,n_1$ even; $n_2$ odd; $n_1\ge2$\\
Gaps&$M\ge6$ even; $5\le N\le2519$ and $\gcd(N,6)=1$\\
Odd part of $M$&At least $5$ (Lemma~\ref{lem:odd-gap-part})\\
$4\mid M$&$m_2\equiv0\pmod4$, $5\nmid N$\\
$M\equiv2\pmod4$&$m_2+n_1\equiv2\pmod4$, $q^N\equiv3\pmod5$\\
Common divisor&$\gcd(g,3570q(q-1))=1$, $\gcd(g,E)=1$, $gE\mid M$\\
$h=M/E$&$h\le630$, $(2h-1)n_1<N$, $h<(N+2)/4$\\
$4\mid r$&$h$ is a fourth power modulo $q$\\
$8\mid r$&$2h$ is an eighth power modulo $q$\\
$5\mid r$&$q\mid F_{(q-1)/5}$ and $h$ is a fifth power modulo $q$\\
$r\equiv4\pmod8$, $h\le32$&$(4h-1)n_1<N$\\
$g>1$&$N\ge5g$, $g\le503$, at most two prime factors with
multiplicity, each at least $11$\\
$N=5g$, $g>1$&$n_1=2$, $M=gE$\\
Guaranteed $g=1$&$N<65$; $3\mid n_2$; any prime $N$; $N=5^a7^b17^j$\\
$N=5d$, $d$ prime&$g=1$ for $d=7,17,29,41,43,73,89,109,113$\\\hline
\end{longtable}}
Section~\ref{sec:forcing-common-factor} gives a way to force $g>1$
using primitive Mersenne factors and their residual indices.

\clearpage
\section{Auxiliary calculations and comparison with MSS}
\label{sec:auxiliary}
% Added in discussion: comparison with MSS and the two adaptations below.

\subsection{Scope of the MSS results}

The system is a special case of $a^x+b^y=c^z$ with
$(a,b,c)=(2,q,p)$ and $z=Z=1$. MSS's complete results for
$c=7,13,97$ exclude those values of $p$ here. Their more general effective
results fix $c$ in specified families, so a bound depending on $c=p$
does not supply a global bound for our problem. Their Pillai theorem
concerns a difference of powers and is not directly a theorem about
our sum. See the theorem statements and outline in \cite{MSS}.

Their polynomial-exponential equation matches our equality exactly:
\[
X^a-X^b=q^{y_1}-q^{y_2},\qquad
(X,a,b,y_1,y_2)=(2,m_1,m_2,n_2,n_1).
\]
But the proposition labeled \texttt{prop-EuAl} in the local MSS source
assumes $(a-b)/E$ odd, where $E$ is the least positive exponent for which
$X^E\equiv\pm1\pmod q$. This is a signed order, unlike our ordinary orders.
For $X=2$ put $r=\operatorname{ord}_q(2)$. If $r$ is odd, $E=r$,
and $M/E$ is even. If $r$ is even, $E=r/2$, and $M/E=2(M/r)$ is even.
Thus the proposition as stated never applies to this specialization.
The polynomial argument below is a separate adaptation, not an invocation
of that proposition.

\subsection{A polynomial identity}
\label{sec:polynomial}

Use the ordinary order and write
\[
r=\operatorname{ord}_q(2),\qquad M=rh,\qquad
\eta=v_q(h),\qquad \rho=v_q(2^r-1).
\]
Here $r\geq2$, $q\nmid r$, $\rho\geq1$, and
$n_1=\rho+\eta$ by lifting the exponent.
Define
\[
F_r(t)=\frac{t^r-1}{t-1},\qquad
J_h(t)=1+t^r+\cdots+t^{r(h-1)},\qquad
G(t)=t^{m_2}(t-1)J_h(t).
\]
We have
\[
F_r(2)=2^r-1,\qquad
F_r(2)G(2)=q^{n_1}(q^N-1).
\]
The polynomial identity below does not require $h$ odd.

Put
\[
H_r(t)=\frac{F_r(t)-r}{t-1}\in\mathbb Z[t],
\qquad c\in\{0,\ldots,r-1\},\quad c\equiv-m_2\pmod r,
\]
and let $R(t)$ be the remainder of $-t^cH_r(t)$ on division by the
monic polynomial $F_r(t)$. Then there is $P(t)\in\mathbb Z[t]$ such that
\begin{equation}\label{eq:adapted-bezout}
\boxed{F_r(t)P(t)+G(t)R(t)=rh=M.}
\end{equation}
To check this, work modulo $F_r(t)$. We have $t^r\equiv1$,
$J_h(t)\equiv h$, and $(t-1)H_r(t)\equiv-r$. Therefore
\[
G(t)R(t)\equiv
-h t^{m_2+c}(t-1)H_r(t)\equiv rh.
\]
Since $F_r$ is monic, the quotient defining $P$ has integer coefficients.

Evaluate \eqref{eq:adapted-bezout} at $2$ and multiply by $F_r(2)$.
Using the original equation gives the exact identity
\[
q^{n_1}R(2)+M(2^r-1)
=(2^r-1)^2P(2)+q^{n_2}R(2).
\]
Consequently
\begin{equation}\label{eq:adapted-congruence}
\boxed{
q^{n_1}R(2)+M(2^r-1)
\equiv0\pmod{q^{\min\{2\rho,n_2\}}}.}
\end{equation}
This is an even-quotient analogue of the congruence used by MSS.

Write $h=q^\eta h_0$ and $2^r-1=q^\rho K$ with $q\nmid h_0K$.
If $\rho>\eta$, division by $q^{n_1}$ gives
\begin{equation}\label{eq:adapted-divided}
\boxed{R(2)+r h_0K\equiv0\pmod{q^{\min\{\rho-\eta,N\}}}.}
\end{equation}
If $\eta\geq\rho$, the modulus in
\eqref{eq:adapted-congruence} is at most $q^{n_1}$, and both left-hand
terms are already divisible by that modulus. It then adds no restriction.
In particular, when $\rho=1$ the calculation gives new restrictions
only if $\eta=0$, and then only modulo $q$.
It is not established whether this improves on direct expansions of the
original equation.

For orientation, $r=2$ gives $q=3$ and $R(t)=1$, since $m_2=1$.
If $\eta=0$, \eqref{eq:adapted-divided} says $1+2h\equiv0\pmod3$,
so $h\equiv1\pmod3$. This illustrates the formula; the $q=3$ branch is already settled by
Lemma~\ref{lem:qthree} without this construction.

To obtain a global estimate, we would need control of the size and sign of $R(2)+r h_0K$, or a
stronger modulus, sufficient to combine with \eqref{eq:baker-M}.
The degree of $R$ can grow with $r$, so a large modulus alone does not
supply that comparison. This is the same kind of size issue emphasized
in MSS's closing discussion.

\subsection{\texorpdfstring{Congruences modulo $p$}{Congruences modulo p}}

Applying the determinant idea in MSS, set
\[
T=m_1n_2-m_2n_1=m_2N+Mn_1+MN>0.
\]
From $2^{m_i}\equiv-q^{n_i}\pmod p$, raising the two congruences
to $n_2,n_1$ respectively and dividing gives
\[
2^T\equiv(-1)^{n_2-n_1}=-1\pmod p.
\]
Similarly, using exponents $m_1,m_2$ gives
\[
q^T\equiv(-1)^{m_1-m_2}=1\pmod p.
\]
Thus
\[
\operatorname{ord}_p(q)\mid T,\qquad
\operatorname{ord}_p(2)\mid2T,\qquad
\operatorname{ord}_p(2)\nmid T.
\]
These are necessary conditions involving the prime $p$ itself.
They have not yet been combined into a new bound.
The normalization of orders used elsewhere in MSS should not be assumed
to preserve our fixed base $2$ without checking the change of variables.

\phantomsection
\addcontentsline{toc}{section}{References}
\bibliographystyle{alpha}
\bibliography{prime_power_notes}
\end{document}
